\exercisesection{2}

\begin{enumerate}
\def\labelenumi{\arabic{enumi}.}
\item
  给出一个单射环同态 \(\rho: \mathbf{A} \to \mathbf{B}\) 的例子，使 \(\mathbf{B}\) 是有限 A-代数，映射 \(^a\rho: \operatorname{Spec}(\mathbf{B}) \to \operatorname{Spec}(\mathbf{A})\) 是满射，但 \(\mathbf{B}\) 不是平坦 A-模。反之，给出一个忠实平坦 A-代数的例子，它是有限生成的，但不在 \(\mathbf{A}\) 上整。
\item
  设 A 是使 \(\text{Spec}(A)\) 为 Hausdorff 的环（第二章 § 4 习题 16）。证明：对每个在 A 上整的 A-代数 B，\(\text{Spec}(B)\) 都是 Hausdorff 的。
\item
  设 A 是环，A' 是包含 A 且在 A 上整的环。证明：对 A 的每个理想 a，a 的根等于 A 与 aA' 的根之交。并证明后者是由满足下述条件的元素 \(x \in A'\) 组成的集合：其整依赖方程的系数（首项系数除外）都属于 a（注意，若 y、z 是 A' 的满足这样一个整依赖方程的两个元素，则 y - z 也满足）。
\item
  设 A 是 Noether 环，\(\rho: A \to B\) 是使 B 成为有限生成 A-模的环同态，N 是有限生成 B-模。证明：诸素理想 \(\mathfrak{p} \in \operatorname{Ass}(\rho_{*}(N))\) 恰是在 \(\rho\) 之下诸素理想 \(\mathfrak{q} \in \operatorname{Ass}(N)\) 的原像。
\item
  设 K 是域；在两个未定元的多项式整环 K{[}X, Y{]} 中考虑子环 A = K{[}X \(^{4}\) , Y \(^{4}\) {]}，A' = K{[}X \(^{4}\) , X \(^{3}\) Y, XY \(^{3}\) , Y \(^{4}\) {]}；A 是 Noether 的，A' 是有限 A-代数。设 p = AY \(^{4}\)，它是 A 的一个素理想；证明 m = A'X \(^{4}\) + A'X \(^{3}\) Y + A'XY \(^{3}\) + A'Y \(^{4}\) 是与 A' 的理想 pA' 相伴的一个（浸入）素理想，但并不位于 p 之上（注意 m 是 X \(^{2}\) Y \(^{6}\) 的剩余类在环 A'/pA' 中的零化子）（参见 § 1 习题 26(d)）。
\item
  定义代数数域的一个递增序列 \(K_{0} = Q, K_{1}, \ldots, K_{n}, \ldots\)，使得：若 \(A_{n}\) 表示 Z 在 \(K_{n}, p\) 中的整闭包，其中 p 是一个素数 \(\neq 2\)，则 \(K_{n}\) 是 \(K_{n-1}\) 的二次扩张，且在 \(A_{n}\) 中存在位于 (p) 之上的一个素理想 \(p_{n}\)，使得在 \(A_{n+1}\) 中有两个不同素理想 \(p_{n+1}, p'_{n+1}\) 位于 \(p_{n}\) 之上（注意，在特征 \(\neq 2\) 的有限域中总存在不是平方的元素）。设 K 是诸 \(K_{n}\) 的并，且 \(A'\) 是 Z 在 K 中的整闭包；证明 \(A'\) 中存在无穷多个位于 Z 的极大理想 (p) 之上的（极大）理想。
\item
  设 A 是 Noether 整环，A' 是它的整闭包。证明：对 A 的每个素理想 p，位于 p 之上的 A' 的素理想只有有限多个。（先化归为 A 是以 p 为极大理想的局部环的情形；考虑它的完备化 \(\hat{A}\) 与 \(\hat{A}\) 的全分式环 B，并注意 A 的分式域 K 可等同于 B 的一个子环（第三章 § 3 第 4 节定理 3 的推论 2）。注意 B 的谱是有限的，因而 B 中只有有限多个两两正交的幂等元。然后设 C \(\supset A\) 是 K 的一个子环，它是有限生成 A-模，从而是半局部的。注意 C 的完备化 \(\hat{C}\) 可等同于 B 的一个子环；最后注意到 \(\hat{C}\) 是有限多个局部环的乘积，其个数等于 C 的极大理想的个数，即可得出结论。）把上述结果推广到 Noether 环在其全分式环中的整闭包（化归为约化环的情形）。
\item
  若局部环 A 的整闭包是局部环，则称 A 为单分支的。
\end{enumerate}

\begin{enumerate}
\def\labelenumi{(\alph{enumi})}
\item
  设 A 是局部整环，K 是它的分式域。为使 A 是单分支的，必须且只需 K 的每个作为有限生成 A-模的子 A-代数都是局部环。
\item
  证明：若 Noether 局部环 A 的完备化是整环，则 A 是单分支的（若 B 是包含于 K 中的有限 A-代数，证明由 B 与 \(\hat{A}\) 生成的 \(\hat{A}\) 的分式域的子环同构于 \(\mathbf{B} \otimes_{\mathbf{A}} \hat{\mathbf{A}}\)，为此利用 A-模 \(\hat{\mathbf{A}}\) 的平坦性以及 \(\mathbf{B}\) 包含于某个含于 \(\mathbf{K}\) 的有限生成自由 A-模之中这一事实；然后利用 (a) 及第三章 § 2 第 13 节命题 19 的推论）。
\end{enumerate}

\begin{enumerate}
\def\labelenumi{\arabic{enumi}.}
\setcounter{enumi}{8}
\tightlist
\item
  设 \(k_0\) 是域，\(A'\) 是无限未定元序列的多项式环 \(k_0[X_n]_{n \in \mathbb{N}}\)；它是一个整闭整环，其分式域为 \(K' = k_0(X_n)_{n \in \mathbb{N}}\)。设 \(\sigma\) 是一个 \(k_0\) -自同构，定义在 \(A'\) 上，满足
\end{enumerate}

\[
\sigma (\mathrm{X} _ {2 n}) = - \mathrm{X} _ {2 n} + \mathrm{X} _ {2 n + 1}, \quad \sigma (\mathrm{X} _ {2 n + 1}) = \mathrm{X} _ {2 n + 1}
\]

其中 \(n \geqslant 0\) 任意；\(\sigma\) 与恒等自同构构成 A' 的一个自同构群 G。设 A 是 G 的不变量环，K 是它的分式域；K' 是 K 的次数为 2 的可分扩张，而 A' 是 A 在 K' 中的整闭包。

\begin{enumerate}
\def\labelenumi{(\alph{enumi})}
\tightlist
\item
  证明 \(\mathbf{A}'\) 是一个具有无限基的 A-模。（若 \(k_0\) 的特征 \(\neq 2\)，证明当记 \(\mathbf{Y}_n = \mathbf{X}_{2n} - \frac{1}{2}\mathbf{X}_{2n+1}\) 时，\(\mathbf{A}\) 是 \(\mathbf{A}'\) 中由诸 \(\mathbf{X}_{2n+1}\) 与诸乘积 \(\mathbf{Y}_n\mathbf{Y}_m\)（\(n \geqslant 0\)，\(m \geqslant 0\)）生成的子环。若 \(k_0\) 的特征为 2，则 \(\mathbf{A}'\) 由诸 \(\mathbf{X}_{2n+1}\) 与
\end{enumerate}

\[
\mathrm{X} _ {2 n} ^ {2} + \mathrm{X} _ {2 n} \mathrm{X} _ {2 n + 1}
\]

（\(n\geqslant 0.\)）生成。

\begin{enumerate}
\def\labelenumi{(\alph{enumi})}
\setcounter{enumi}{1}
\item
  设 \(\mathfrak{m}'\) 是 \(\mathbf{A}'\) 的由所有 \(\mathbf{X}_n\) 生成的极大理想，并设 \(\mathfrak{m} = \mathbf{A} \cap \mathfrak{m}'\)；证明 \(\mathfrak{m}' = \mathfrak{mA}'\)，\(\mathfrak{m}'\) 是 \(\mathbf{A}'\) 中位于 \(\mathfrak{m}\) 之上的唯一理想，且域 \(\mathbf{A} / \mathfrak{m}\) 与 \(\mathbf{A}' / \mathfrak{m}'\) 典范同构。
\item
  设 \(\mathfrak{p}'\) 是 \(\mathbf{A}'\) 的由诸 \(\mathbf{X}_{2n+1}\)（\(n \geqslant 0\)）生成的素理想，并设 \(\mathfrak{p} = \mathfrak{p}' \cap \mathbf{A}\)；设 \(k, k'\) 分别是 \(\mathbf{A}/\mathfrak{p}\) 与 \(\mathbf{A}'/\mathfrak{p}'\) 的分式域；\(\mathfrak{p}'\) 是 \(\mathbf{A}'\) 中位于 \(\mathfrak{p}\) 之上的唯一素理想。若 \(k_0\) 的特征不是 2，则 \(k'\) 是 \(k\) 的次数为 2 的可分扩张；若 \(k_0\) 的特征为 2，则 \(k'\) 是 \(k\) 的无限根式扩张。
\end{enumerate}

\begin{enumerate}
\def\labelenumi{\arabic{enumi}.}
\setcounter{enumi}{9}
\tightlist
\item
  \begin{enumerate}
  \def\labelenumii{(\alph{enumii})}
  \tightlist
  \item
    设 \(\mathbf{A}'\) 是两个未定元的多项式环 \(\mathbf{Z}[\mathbf{X},\mathbf{Y}]\)，\(\sigma\) 是 \(\mathbf{A}'\) 的保持 \(\mathbf{Y}\) 不变且满足 \(\sigma (\mathbf{X}) = -\mathbf{X}\) 的自同构；\(\sigma\) 与恒等自同构构成一个自同构群 \(\mathcal{G}\)，作用于 \(\mathbf{A}'\)，\(\mathcal{G}\) 的不变量环为 \(\mathbf{A} = \mathbf{Z}[\mathbf{X}^2,\mathbf{Y}]\)。设 \(\mathfrak{p}'\) 为素理想 \(\mathrm{A}'(\mathrm{Y} - 2\mathrm{X})\)，\(\mathfrak{p} = \mathrm{A}\cap \mathfrak{p}'\)；分解群 \(\mathcal{G}^{\mathbb{Z}}(\mathfrak{p}')\) 即单位元。证明 \(\mathrm{A / p}\) 与 \(\mathrm{A}' / \mathfrak{p}'\) 不同构（考虑这两个环与 \(\mathbf{Z} / 2\mathbf{Z}\) 的张量积）。
  \end{enumerate}
\end{enumerate}

\begin{enumerate}
\def\labelenumi{(\alph{enumi})}
\setcounter{enumi}{1}
\tightlist
\item
  设 K 是特征 \(\neq2\) 的域，A' 是多项式环 K{[}X, Y{]}，G 是 A' 的自同构群，由恒等映射与 A' 的 K{[}Y{]}-自同构 \(\sigma\)（满足 \(\sigma(\mathbf{X}) = -\mathbf{X}\)）组成；G 的不变量环 A 为 \(K[X^{2}, Y]\)。设 \(p'\) 是 A' 的素理想 \(\mathbf{A}'(\mathbf{X}^{3} - \mathbf{Y})\)，\(p = A \cap p'\)；分解群 \(\mathcal{G}^{\mathrm{Z}}(p')\) 即单位元，但 A/p 与 \(A'/p'\) 不同构。
\end{enumerate}

¶ 11. (a) 设 A 是局部整闭整环，m 是它的极大理想，K 是它的分式域，K' 是 K 的有限 Galois 扩张，A' 是 A 在 K' 中的整闭包，m' 是 A' 的位于 m 之上的一个极大理想，B = \(\mathbf{A}^{\mathbb{Z}}(\mathfrak{m}')\) 是它的分解环，n = m' ∩ B。对每个整数 k \textgreater{} 0，有 B = A + n\^{}k（按命题 4 的方式进行论证）。此外，存在 \(\mathbf{K}^{\mathbf{Z}}\)（B 的分式域）的一个本原元 z 使得 z ∈ n；由此推出 n ∩ A{[}z{]} = mA{[}z{]}。

\begin{enumerate}
\def\labelenumi{(\alph{enumi})}
\setcounter{enumi}{1}
\item
  证明：对每个整数 \(k\)，有 \(\mathfrak{m}^k\mathrm{B}_{\mathfrak{n}} \cap \mathrm{A} = \mathfrak{m}^k\)。（若 \(x \in \mathfrak{m}^k\mathrm{B}_{\mathfrak{n}} \cap \mathrm{A}\)，注意存在 \(t \in \mathrm{A}[z]\)（其中 \(z\) 如 (a) 中所定义），使得 \(t \notin \mathfrak{n}\) 且 \(tx \in \mathfrak{m}^k\mathrm{A}[z]\)；记 \(t = t_0 + t_1z + \dots + t_{r-1}z^{r-1}\)（其中 \(r\) 是次数 \([\mathbf{K}^{\mathbf{Z}}: \mathbf{K}]\)），由此推出 \(t_0x \in \mathfrak{m}^k\) 且 \(t_0 \notin \mathfrak{m}\)。）
\item
  现在设 \(\mathbf{K}'\) 是 \(\mathbf{K}\) 的任意 Galois 扩张。保持与上面相同的记号，证明仍有 \(\mathbf{B} = \mathbf{A} + \mathfrak{n}^k\) 与 \(\mathfrak{m}^k\mathrm{B}_{\mathfrak{n}} \cap \mathbf{A} = \mathfrak{m}^k\)（注意 \(\mathbf{A}^{\mathbb{Z}}(\mathfrak{m}')\) 的每个元素都包含在 \(\mathfrak{m}' \cap \mathbf{C}\) 的分解环中，其中 \(\mathbf{C}\) 是 \(\mathbf{A}\) 在 \(\mathbf{K}'\) 的一个有限 Galois 子扩张中的整闭包）。
\item
  在 (c) 的假设下，证明：若 A 是 Noether 的，则 \(B_{n}\) 亦然。（注意 \(B_{n}\) 关于 m-进拓扑的 Hausdorff 完备化依 (c) 可等同于 A 的完备化 \(\hat{A}\)；若 \(\phi: B_{n} \rightarrow \hat{A}\) 是
\end{enumerate}

典范同态，证明：对每个理想 \(a\)（\(B_n\) 中的任意理想），有 \(\phi(a\hat{A}) = a\)，为此利用 \(\hat{A}\) 中每个理想都是有限生成的这一事实，从而 \(a\hat{A}\) 由 \(a\) 的有限多个元素生成，于是化归为 \(K'\) 是 \(K\) 的有限扩张的情形。）

\begin{enumerate}
\def\labelenumi{\arabic{enumi}.}
\setcounter{enumi}{11}
\tightlist
\item
  设 K 是域，\(K'\) 是 K 的有限 Galois 扩张，C 是多项式环 \(K'[X, Y]\)，B 是商环 C/CXY，x、y 是 X、Y 在 B 中的典范像。设 \(A'\) 是 B 的子环 \(K[x] + yK'[y]\)。\(K'\) 在 K 上的 Galois 群 G 忠实作用于 \(A'\)，G 在 \(A'\) 中的不变量环 A 为 \(K[x, y]\)。设 S 是 \(A'\) 中 x 的幂组成的集。证明 \(S^{-1}A' = S^{-1}A\)，且 G 并不忠实作用于 \(S^{-1}A'\)。
\end{enumerate}

¶ 13. (a) 设 K 是特征 0 的域，n 是多项式环 K{[}X, Y, Z{]} 的由 \(Y^{2} - X^{2} - X^{3}\) 生成的理想。证明 n 是素的，且整环 A = K{[}X, Y, Z{]}/n 不是整闭的：若 x、y、z 是 X、Y、Z 在 A 中的像，则 A 的分式域 L 中的元素 t = y/x 在 A 上整，但不属于 A。设 \(A' \supset A\) 是 L 的子环 K{[}x, t, z{]}，它在 A 上整。在 A 中考虑由 xz - y 与 \(z^{2} - 1 - x\) 生成的素理想 p，以及由 x、y、z - 1 生成的极大理想 \(q \supset p\)。另一方面，在 A' 中考虑极大理想 \(q'\)，它由 x、\(t + 1\) 与 z - 1 生成；证明 \(q'\) 位于 q 之上，但不存在 A' 的位于 p 之上的素理想 \(p' \subset q'\)。（考虑一个从 \(A'/pA'\) 到 K 的扩域的同态，并证明在这样的同态下 \(q'\) 的像必为 0。）

\begin{enumerate}
\def\labelenumi{(\alph{enumi})}
\setcounter{enumi}{1}
\item
  设 \(\mathbf{A}'\) 是多项式环 \(\mathbf{Z}[\mathbf{X}]\) 对理想 \(n\)（由 2X 与 \(\mathbf{X}^2 -\mathbf{X}\) 生成）的商环；设 \(x\) 是 \(\mathbf{X}\) 在 \(\mathbf{A}'\) 中的像；\(\mathbf{A}'\) 在 \(\mathbf{Z}\) 上整，但数 2 是 \(\mathbf{A}'\) 中的零因子。设 \(q'\) 是 \(\mathbf{A}'\) 的由 2 与 \(x - 1\) 生成的素理想；则 \(q' \cap \mathbf{Z} = 2\mathbf{Z}\)；若 \(p\) 是 \(\mathbf{Z}\) 中的素理想 (0)，证明不存在素理想 \(p' \subset q'\)（含于 \(\mathbf{A}'\) 且位于 \(p\) 之上）。
\item
  设 K 是域，n 是多项式环 K{[}X, Y, Z{]} 的由 \(X^{2}\)、XY、XZ、YZ - Y 与 \(Z^{2} - Z\) 生成的理想，A' 是商环 K{[}X, Y, Z{]}/n，x、y、z 是 X、Y、Z 在 A' 中的像，A 是 A' 的子环 K{[}x, y{]}；A 在它的全分式环中整闭，而 A' 在 A 上整。在 A 中考虑由 x 生成的素理想 p 以及由 x 与 y 生成的极大理想 q。另一方面，在 A' 中考虑由 x、y、z 生成的极大理想 \(q'\)；证明 \(q'\) 位于 q 之上，但不存在位于 p 之上的素理想 \(p' \subset q'\)（方法同 (a)）。
\end{enumerate}

* 试从几何上解释 (a) 与 (c) 中的例子。 *

¶ 14. (a) 设 A 是域 K 的子环，x 是 K 中一个 \(\neq0\) 的元素。证明：若 x 不在 A 上整，则存在 \(A[x^{-1}]\) 的一个包含 \(x^{-1}\) 的极大理想，且对 \(A[x^{-1}]\) 的每个包含 x 的极大理想 m，理想 \(A\cap m\) 都是 A 的极大理想（注意 x 不属于环 \(A[x^{-1}]\)）。

\begin{enumerate}
\def\labelenumi{(\alph{enumi})}
\setcounter{enumi}{1}
\item
  设 A 是局部整环，m 是它的极大理想，K 是包含 A 的域，k 是 A 的剩余域。设 \(x \in K\) 满足 \(x \notin A\) 且 \(x^{-1} \notin A\)，并设 A 在环 A{[}x{]} 中整闭。证明存在一个 A-代数同态 \(A[x] \to k[X]\)，它延拓典范同态 \(A \to k\)，且把 x 映为 X。（需要证明：若 \(\sum_{i=0}^{n}a_{i}x^{i}=0\)，其中对所有 i 有 \(a_{i}\in A\)，则所有 \(a_{i}\) 都属于 m。否则将有一个关系式 \(a_{n}x^{n-p}+\cdots+a_{p}=-(a_{p+1}x^{-1}+\cdots+a_{0}x^{-p})\)，其中 \(a_{p}\) 在 A 中可逆；利用 §1 习题 5，证明两端的公共值 b 属于 A，再利用 (a)，证明必有 \(b\in m\)；由此推出 \(x^{-1}\) 将在 A 上整，这是荒谬的。）
\item
  在 (b) 的假设下，证明 A{[}x{]} 的理想 mA{[}x{]} 是素的且非极大的。
\end{enumerate}

\begin{enumerate}
\def\labelenumi{\arabic{enumi}.}
\setcounter{enumi}{14}
\item
  设 A 是环，p 是 A 的素理想。证明：在多项式环 B = A{[}X{]} 中，理想 pB = q 是素的；局部环 \(B_{q}\) 的极大理想等于 \(pB_{q}\)；若 k 是 A/p 的分式域，则 \(B_{q}/pB_{q}\) 同构于有理分式域 \(k(\mathbf{X})\)。若 A 是整闭的，则 \(B_{q}\) 亦然。
\item
  设 A 是整环，K 是它的分式域，\(K'\) 是 K 的次数为 n 的有限扩张，\(A'\) 是 \(K'\) 中包含 A 的子环，以 \(K'\) 为分式域且在 A 上整。设 m 是 A 的一个极大理想，B 是局部环 \(A[X]_{mA[X]}\)（习题 15）。
\end{enumerate}

\begin{enumerate}
\def\labelenumi{(\alph{enumi})}
\item
  子环 \(\mathbf{B}' = \mathbf{B}[\mathbf{A}']\) 含于 \(\mathbf{K}'(\mathbf{X})\)，它在 \(\mathbf{B}\) 上整；它的分式域是 \(\mathbf{K}'(\mathbf{X})\)，且 \([\mathbf{K}'(\mathbf{X}) : \mathbf{K}(\mathbf{X})] = n\)。
\item
  设 \(\mathfrak{m}'\) 是 \(\mathbf{A}'\) 的包含 \(\mathfrak{m}\) 的任一极大理想，则 \(\mathfrak{n}' = \mathfrak{m}'\mathbf{B}'\) 是 \(\mathbf{B}'\) 的极大理想；若记 \(\mathfrak{n} = \mathfrak{m}\mathbf{B}\)，则
\end{enumerate}

\[
[ (\mathbf {A} ^ {\prime} / \mathfrak {m} ^ {\prime}): (\mathbf {A} / \mathfrak {m}) ] = [ (\mathbf {B} ^ {\prime} / \mathfrak {n} ^ {\prime}): (\mathbf {B} / \mathfrak {n}) ]
\]

以及

\[
[ (\mathrm{A} ^ {\prime} / \mathfrak {m} ^ {\prime}): (\mathrm{A} / \mathfrak {m}) ] _ {s} = [ (\mathrm{B} ^ {\prime} / \mathfrak {n} ^ {\prime}): (\mathrm{B} / \mathfrak {n}) ] _ {s}.
\]

（注意 \((\mathbf{A}' / \mathfrak{m}') \otimes_{\mathbf{A}} \mathbf{B}\) 同构于域 \((\mathbf{A}' / \mathfrak{m}')(\mathbf{X})\)，由此推出 \(\mathbf{B}' / \mathfrak{n}'\) 同构于 \((\mathbf{A}' / \mathfrak{m}')(\mathbf{X}).\)）

\begin{enumerate}
\def\labelenumi{(\alph{enumi})}
\setcounter{enumi}{2}
\tightlist
\item
  映射 \(m' \mapsto m'B'\) 是从 \(A'\) 的包含 m 的极大理想所成之集到 \(B'\) 的极大理想所成之集上的双射；逆双射是 \(n' \mapsto n' \cap A'\)。
\end{enumerate}

¶ 17. (a) 设 K 是无限域，E 是 K 的有限可分代数扩张，g 是 K{[}X{]} 的一个多项式。证明存在元素 \(x \in E\)，使得 \(\mathrm{E} = \mathrm{K}(x)\)，且 x 在 K 上的极小多项式与 g 互素（按《代数》第五章 § 7 第 7 节命题 12 的方式论证）。

\begin{enumerate}
\def\labelenumi{(\alph{enumi})}
\setcounter{enumi}{1}
\item
  设 \(k\) 是域，A 是准素 \(k\) -代数（第四章 § 2 习题 32），\(m\) 是它的唯一素理想。若 \(A / m\) 是 \(k\) 的超越扩张，则 A 中存在不在 \(k\) 上代数的元素。若 \(A / m\) 在 \(k\) 上代数，且 \(d\) 是一个至少等于 \(A / m\) 在 \(k\) 上的次数的可分部分的整数（从而当该可分部分为无限时，d 可为任意整数），证明 A 中存在在 \(k\) 上的极小多项式的次数 \(\geqslant d\) 的元素。进而，若 \(A / m\) 在 \(k\) 上不可分，或 \(m \neq 0\)，则 A 中存在一个元素，它在 \(k\) 上的极小多项式的次数 \(>d\)。（若 \(A / m\) 在 \(k\) 上不可分，利用《代数》第五章 § 8 习题 5。若 \(A / m\) 的次数为 \(d\) 且在 \(k\) 上可分，又 \(x \in A\) 使 \(\xi \in A / m\)（即 \(x\) 的剩余类）是 \(A / m\) 的本原元，其极小多项式为 \(f \in k[X]\)，则注意 \(f'(\xi) \neq 0\)，并由此推出：若 \(y \neq 0\) 是 \(m\) 的一个元素，则 \(f(x + y) \neq 0\)。）
\item
  设 \(k\) 是无限域，\(A\) 是 \(k\) -代数，为 \(r\) 个准素代数 \(A_{i} (1 \leqslant i \leqslant r)\) 的直合成；设 \(m_{i}\) 是 \(A_{i}\) 的极大理想，\(d_{i}\) 是一个整数，它至少等于 \(A_{i} / m_{i}\) 在 \(k\) 上的次数的可分部分（当 \(A_{i}\) 代数且该部分有限时），否则为任意整数。证明 \(A\) 中存在一个不在 \(k\) 上代数的元素，或者
\end{enumerate}

其在 \(k\) 上的极小多项式的次数 \(\geqslant \sum_{i=1}^{r} d_i\)；进而，若某个 \(m_i\) 为 \(\neq 0\)，或某个域 \(A_i / m_i\) 不是 \(k\) 的可分代数扩张，则 \(A\) 中存在一个不在 \(k\) 上代数的元素，或者其极小

多项式在 \(k\) 上的次数 \(> \sum_{i=1}^{r} d_i\)（利用 (a) 与 (b)）。

¶ 18. 设 A 是局部整闭整环，K 是它的分式域，m 是它的极大理想，\(k = \mathbf{A} / \mathfrak{m}\) 是它的剩余域。设 \(\mathbf{K}'\) 是 \(\mathbf{K}\) 的次数为 \(n\) 的有限代数扩张，\(\mathbf{A}'\) 是 \(\mathbf{K}'\) 中包含 \(\mathbf{A}\) 的子环，它在 \(\mathbf{A}\) 上整且以 \(\mathbf{K}'\) 为分式域，\(\mathfrak{m}_i\)（\(1 \leqslant i \leqslant r\)）是 \(\mathbf{A}'\) 的诸极大理想。称理想 \(\mathfrak{m}\) 在 \(\mathbf{A}'\) 中是非分歧的，如果存在 \((w_i)_{1 \leqslant i \leqslant n}\)，它构成 \(\mathbf{K}'\) 在 \(\mathbf{K}\) 上的一组基，由 \(\mathbf{A}'\) 的元素组成，且其判别式 \(D_{\mathbf{K}' / \mathbf{K}}(w_1, \ldots, w_n)\)（《代数》第九章 §2）属于 \(\mathbf{A} - \mathfrak{m}\)。

\begin{enumerate}
\def\labelenumi{(\alph{enumi})}
\item
  证明：若 \(\mathfrak{m}\) 在 \(\mathbf{A}'\) 中非分歧，则 \(\mathbf{K}'\) 是 \(\mathbf{K}\) 的可分扩张，\(\mathbf{A}'\) 是 \(\mathbf{A}\) 在 \(\mathbf{K}'\) 中的整闭包，并且是自由 \(\mathbf{A}\) -模，它在 \(\mathbf{A}\) 上的任何基都是 \(\mathbf{K}'\) 在 \(\mathbf{K}\) 上的基；\(\mathfrak{mA}'\) 是 \(\mathbf{A}'\) 的 Jacobson 根，而 \(\mathbf{A}' / \mathfrak{mA}'\) 是 \(k\) 上秩为 \(n\) 的半单 \(k\) -代数，它在 \(k\) 上可分，且是诸域 \(\mathbf{A}' / \mathfrak{m}_i'\) 的直合成。
\item
  设 \(d_{i}\) 是 \(A'/m_{i}'\) 在 \(A/m\) 上的次数的可分部分。证明：若 \(k\) 是无限的，则 \(\sum_{i=1}^{r} d_{i} \leqslant n\)（利用习题 17 (c)）；此外，下列条件等价：
\end{enumerate}

（α）m 在 A' 中非分歧。

\((\beta)\) \(\mathbf{A}'\) 是有限生成 A-模，且 \(\mathbf{A}' / \mathfrak{mA}'\) 是半单 \(k\) -代数并在 \(k\) 上可分。

\((\gamma)\sum_{i=1}^{r}d_{i}\leqslant n.\)

（为看出 \((\beta)\) 蕴涵 \((\gamma)\)，注意，利用第二章 § 3 第 2 节命题 4 的推论 2，A-模 \(A'\) 有一个生成系，其元素模 \(mA'\) 的剩余类构成 \(A'/mA'\) 在 \(k\) 上的一组基；由此推出

\[
[ \mathbf {K} ^ {\prime}: \mathbf {K} ] \leqslant [ (\mathbf {A} ^ {\prime} / \mathsf {m A} ^ {\prime}): k ].
\]

为看出 \((\gamma)\) 蕴涵 \((\alpha)\)，利用习题 17 (c) 证明：存在 \(\mathbf{A}' / \mathfrak{mA}'\) 在 \(k\) 上的一组基，由一个元素的诸幂 \(\xi^i\)（\(\xi\) 为该元素，\(0 \leqslant i \leqslant n - 1\)）组成；若 \(x \in \mathbf{A}'\) 是剩余类 \(\xi\) 的一个元素，则推出 \(\mathrm{D}_{\mathrm{K}' / \mathrm{K}}(1, x, \ldots, x^{n-1})\) 属于 \(\mathbf{A} - \mathfrak{m}\)。）给出一个例子，其中 \(\mathbf{K}'\) 在 \(\mathbf{K}, \mathbf{A}' / \mathfrak{mA}' = \mathbf{A} / \mathfrak{m}\) 上可分，而 \(\mathbf{A}'\) 不是有限生成 A-模（参见习题 9 (c)）。

\begin{enumerate}
\def\labelenumi{(\alph{enumi})}
\setcounter{enumi}{2}
\tightlist
\item
  把 (b) 的结果推广到 \(k\) 为有限的情形（在习题 16 的记号下，证明：若 \(n\) 在 \(B'\) 中非分歧，则 \(m\) 在 \(A'\) 中非分歧；利用 \(B / n\) 是无限的这一事实，使用习题 16 (c)，从而得到一个元素 \(z \in B'\)，使得
\end{enumerate}

\[
\mathrm{D} _ {\mathbf {K} ^ {\prime} (\mathbf {X}) / \mathbf {K} (\mathbf {X})} (1, z, \dots , z ^ {n - 1}) \in \mathrm{B} - n;
\]

证明可设 \(z\) 是 \(\mathbf{A}'[\mathbf{X}]\) 中的一个多项式，并且把上面的判别式表为 \(\mathbf{X}\) 的多项式，便得到由 \(n\) 个元素 \(w_{i}\)（\(1 \leqslant i \leqslant n\)）组成的一个组，属于 \(\mathbf{A}'\)，使得 \(\mathrm{D}_{\mathrm{K}'/\mathrm{K}}(w_1, \ldots, w_n) \in \mathbf{A} - \mathfrak{m}\)）。

¶ 19. 设 A 是整闭整环，K 是它的分式域，K' 是 K 的次数为 n 的有限代数扩张，A' 是 K' 中包含 A 的子环，它在 A 上整且以 K' 为分式域；若 p 是 A 的素理想，用 \(\mathbf{A}_{\mathfrak{p}}^{\prime}\) 表示 \(\mathbf{A}^{\prime}\) 的分母属于 \(\mathbf{A}-\mathfrak{p},\mathfrak{p}\) 的分式环；则称 p 在 \(\mathbf{A}^{\prime}\) 中非分歧，如果 \(\mathfrak{p}\mathbf{A}_{\mathfrak{p}}\) 在 \(\mathbf{A}_{\mathfrak{p}}^{\prime}\) 中非分歧。

\begin{enumerate}
\def\labelenumi{(\alph{enumi})}
\tightlist
\item
  设 \(\mathfrak{p}_i^{\prime}(1\leqslant i\leqslant r)\) 是 \(\mathbf{A}'\) 的位于 \(\mathfrak{p}\) 之上的诸素理想。设 \(d_{i}\) 是 \(\mathbf{A}' / \mathfrak{p}_i'\) 的分式域在
\end{enumerate}

A/p 的分式域上的次数的可分部分；证明 \(\sum_{i=1}^{r} d_i \leqslant n\)。（化归为习题 18。）

\begin{enumerate}
\def\labelenumi{(\alph{enumi})}
\setcounter{enumi}{1}
\tightlist
\item
  下列条件等价：
\end{enumerate}

（α）理想 p 在 A' 中非分歧。

（β）A' 包含 K' 在 K 上的一组基，其判别式属于 A - p。

\((\gamma)\sum_{i=1}^{r}d_{i}=n.\)

特别地，若 \(r = n\)，则 \(\mathfrak{p}\) 在 \(\mathbf{A}'\) 中非分歧。

给出一个例子，其中 \(p\) 在 \(A'\) 中非分歧，但 \(A'\) 不是有限生成 \(A\) -模（参见习题 9 (a)）。

\begin{enumerate}
\def\labelenumi{(\alph{enumi})}
\setcounter{enumi}{2}
\item
  若 A 的每个素理想都在 A' 中非分歧，则称 A 在 A'（或在 K'）中非分歧。证明此时 A' 是忠实平坦 A-模（参见第二章 § 3 第 4 节命题 15 的推论）。设 K'' 是 K' 的次数为 n' 的有限代数扩张；设 A' 是 A 在 K' 中的整闭包，A'' 是 A 在 K'' 中的整闭包；为使 A 在 A'' 中非分歧，必须且只需 A 在 A' 中非分歧且 A' 在 A'' 中非分歧。
\item
  设 \(k\) 是域，\(A\) 是整闭的整 \(k\) -代数，\(K\) 是它的分式域。设 \(k'\) 是 \(k\) 的有限可分代数扩张，次数为 \(n\)（在 \(k\) 上）。证明：若 \(k' \otimes_k K\) 是域，则 \(A\) 在 \(k' \otimes_k K\) 中非分歧（考虑本原元 \(x\)（\(k'\) 在 \(k\) 上的）以及 \(k' \otimes_k K\) 的由诸 \(x^i\)（\(0 \leqslant i \leqslant n - 1\)）组成的基）。
\end{enumerate}

\begin{enumerate}
\def\labelenumi{\arabic{enumi}.}
\setcounter{enumi}{19}
\item
  设 A 是第三章 § 3 习题 15 (b) 中所定义的整环，C 是局部环 \(\mathbf{A}_{\mathfrak{m}}\)，\(\mathfrak{n} = \mathfrak{mA}_{\mathfrak{m}}\) 是 C 的极大理想，L 是 A 与 C 的分式域。若 \(x, y\) 是 X 与 Y 在 C 中的典范像，证明 \(t = y / x\) 不属于 C，但 \(t^2 \in \mathbf{C}\)，从而 C 不是整闭的；此外 \(\mathbf{C}' = \mathbf{C}[t]\) 是 C 的整闭包，且同构于 \(\mathbf{S}^{-1}\mathbf{K}[\mathbf{T}]\)（T 为未定元），其中 S 是 K{[}T{]} 中由既不被 T - 1 整除、也不被 T + 1 整除的多项式组成的乘性子集。证明 \(\mathbf{C}'\) 是有限生成 C-模，且 \(\mathbf{C}'\) 的位于 \(\mathfrak{n}\) 之上的（极大）理想是主理想 \(\mathfrak{q}_1 = (t - 1)\mathbf{C}'\) 与 \(\mathfrak{q}_2 = (t + 1)\mathbf{C}'\)；但局部环 \(\mathbf{C}_{\mathfrak{q}_1}'\) 与 \(\mathbf{C}_{\mathfrak{q}_2}'\) 不是在 C 上整的 C-代数。
\item
  设 A 是包含域 Q 的整闭整环，K 是它的分式域，L 是 K 的代数扩张，B 是 A 在 L 中的整闭包。证明：对 A 的每个理想 a，有 \(A \cap aB = a\)。
\item
  设 A 是 Noether 整环，K 是它的分式域，A' 是 K 中包含 A 的子环。设 A 在 A' 中整闭，且对 A 的每个素理想 p，存在素理想 \(p'\)（属于 \(A'\)），位于 p 之上。用反证法论证：设 \(x \in A' \cap C A\)，令 \(a = A \cap (x^{-1} A)\)，它是异于 A 的理想。设 p 是与 A/a 相伴的素理想；则传递子 b = a: p 异于 a（第四章 §1 第 4 节命题 9），因而存在 \(y \in b\)，使得 \(xy \notin A\)，且
\end{enumerate}

\[
x y p \subset p A ^ {\prime} \cap A = p;
\]

由此得出矛盾。）

\exercisesection{3}

¶ 1. 设 A 是环，p 是 A 的素理想，B 是有限生成 A-代数，并设 \(r_{1} \subset \cdots \subset r_{m}\) 是 B 的位于 p 之上的非空递增理想序列。证明存在元素 \(a \in A - p\) 与 B 的元素的一个有限族 \((y_{k})_{1 \leqslant k \leqslant q}\)，具有下列性质：

\begin{enumerate}
\def\labelenumi{(\arabic{enumi})}
\item
  若记 \(\mathbf{B}' = \mathbf{A}[y_1, \ldots, y_q]\)，且 \(\mathbf{S}\) 是 \(a\) 的幂所成之集，则环 \(\mathbf{S}^{-1}\mathbf{B}\) 在 \(\mathbf{S}^{-1}\mathbf{B}'\) 上整。
\item
  对每个 \(i\)（\(1 \leqslant i \leqslant m\)），存在整数 \(h(i) \geqslant 0\)，使得 \(\mathfrak{r}_i\) 包含 \(y_1, \ldots, y_{h(i)}\)，且对每个多项式 \(\mathbf{P} \in \mathbf{A}[\mathbf{X}_{h(i)+1}, \ldots, \mathbf{X}_q]\)，关系 \(\mathrm{P}(y_{h(i)+1}, \ldots, y_q) \in \mathfrak{r}_i\) 蕴涵 \(\mathbf{P}\) 的所有系数都属于 \(\mathfrak{p}\)。
\end{enumerate}

此时理想 \(\mathfrak{r}_i\cap \mathbf{B}'\) 由 \(\mathfrak{p}\) 与诸 \(y_{k}\)（其中 \(k\leqslant h(i)\)）生成。（对 \(m\) 用归纳法，从而化归为 \(m = 1\) 的情形。然后设 \(\mathfrak{r} = \mathfrak{r}_1\alpha\)；对 A-代数 B 的生成元个数 \(n\) 用归纳法。设 \((x_{i})_{1\leqslant i\leqslant n}\) 是这样一个生成系，\(\mathfrak{s}\) 是 \(\mathrm{A}[\mathrm{X}_1,\dots ,\mathrm{X}_n]\) 中由满足 \(\mathrm{P}(x_1,\ldots ,x_n)\in \mathfrak{r}\) 的那些 P 组成的理想；则 \(\mathfrak{s}\supset \mathfrak{pA}[\mathrm{X}_1,\ldots ,\mathrm{X}_n]\)；按 \(\mathfrak{s}\) 是否等于 \(\mathfrak{pA}[\mathrm{X}_1,\ldots ,\mathrm{X}_n]\) 区分两种情形。在第二种情形，存在多项式 \(\mathrm{Q}\in \mathrm{A}[\mathrm{X}_1,\ldots ,\mathrm{X}_n]\)，它的所有系数都不在 \(\mathfrak{p}\) 中，且它属于 \(\mathfrak{s}\)。证明可以确定整数 \(m_i\)（\(2\leqslant i\leqslant n\)），使得当记 \(x_1' = \mathrm{Q}(x_1,\ldots ,x_n),x_i' = x_i - x_1^{m_i}\)（\(2\leqslant i\leqslant n\)）时，存在元素 \(b\in \mathrm{A} - \mathfrak{p}\)，使 \(bx_{1}\) 在 \(\mathrm{C} = \mathrm{A}[x_2',\ldots ,x_n']\) 上整；把归纳假设应用于 A-代数

\[
\mathrm{B} _ {1} = \mathrm{A} [ x _ {2} ^ {\prime}, \dots , x _ {n} ^ {\prime} ]
\]

以及理想 \(\mathfrak{r} \cap \mathbf{B}_1\)。）

\begin{enumerate}
\def\labelenumi{\arabic{enumi}.}
\setcounter{enumi}{1}
\item
  设 K 是域，B 是两个未定元的多项式环 K{[}X, Y{]} 对 K{[}X, Y{]} 的由 \(X^{2}\) 与 XY 生成的理想 a 的商环；用 x、y 表示元素 X、Y 在 B 中的像；设 A 是 B 的子环 K{[}x{]}。证明元素 x 与 \(y^{n}\)（\(n \geqslant 0\)）构成 B 在 K 上的一组基，且 B 不在 A 上整。另一方面，B 中不存在在 A 上代数自由的元素。由此给出当 A 不是整环时第 1 节定理 1 的推论 1 的一个反例。
\item
  设 K 是无限域，L 是 K 的扩域，\((x_{i})_{1 \leq i \leq n}\) 是 L 中元素的有限族，d 是 \(\mathrm{K}(x_{1}, \ldots, x_{n})\) 在 K 上的超越次数。证明存在 L 的 d 个元素 \(z_{j} = \sum_{i=1}^{n} a_{ji} x_{i}\)（其中 \(a_{ji} \in K\)），使得 \(\mathrm{K}[x_{1}, \ldots, x_{n}]\) 在 \(\mathrm{K}[z_{1}, \ldots, z_{d}]\) 上整。若进一步设 \(\mathrm{K}(x_{1}, \ldots, x_{n})\) 是 K 的可分扩张，证明诸 \(a_{ji}\) 可如此选取，使诸 \(z_{j}\) 构成 \(\mathrm{K}(x_{1}, \ldots, x_{n})\) 在 K 上的一个分离超越基（当 n \textgreater{} d 时对 n 用归纳法）。
\item
  设 B 是整环，A 是 B 的子环。设存在 B 的 \(n \geqslant 1\) 个元素 \(t_{1}, \ldots, t_{n}\)，它们在 A 上代数无关，且 B 在 \(A[t_{1}, \ldots, t_{n}]\) 上整。设 n 是 B 的一个极大理想，\(p = n \cap A\)。则存在 B 的素理想 q，它含于 n、异于 n 且包含 p。（先通过考虑 A 的整闭包 \(A'\)（含于 B 的分式域中）与环 \(B' = B[A']\)，化归为 A 整闭的情形；在 A 整闭的情形，利用 §2 第 4 节定理 3。）
\end{enumerate}

¶ 5. 设 A 是整环，B ⊃ A 是 A 的分式域 K 的子环；设 A 在 B 中整闭，且 B 包含元素 t，使 B 是有限生成 A{[}t{]}-模。

\begin{enumerate}
\def\labelenumi{(\alph{enumi})}
\tightlist
\item
  设 F 是一个 \(\neq 0\) 的多项式，属于 A{[}X{]}，使 \(F(t)\) 属于导子 \(\mathfrak{f}\)（B 相对于 A{[}t{]} 的）；证明：若 \(c\) 是 F 的首项系数，则 \(c\) 属于 \(\mathfrak{f}\) 的根。（通过考虑环 A{[} \(c^{-1}\) {]} 与 B{[} \(c^{-1}\) {]}，化归为证明：当 \(c = 1\) 时必有 B = A{[}t{]}。另一方面，可设 F(t) \(\neq 0\)，否则 \(t \in \mathbf{A}\)。对每个 \(x \in \mathbf{B}\)，由假设
\end{enumerate}

\[
x \mathrm{F} (t) = \mathrm{F} (t) \mathrm{Q} (t) + \mathrm{R} (t)
\]

其中 Q 与 R 是 A{[}X{]} 中的多项式，且 \(\deg(\mathrm{R}) < \deg(\mathrm{F})\)；只需考虑 \(\mathrm{R}(t) \neq 0\) 的情形；若 \(y = \mathrm{R}(t)(\mathrm{F}(t))^{-1} = x - \mathrm{Q}(t)\)，证明 t 在 \(A[y^{-1}]\) 上整，并由此推出 y 在 A 上整。）

\begin{enumerate}
\def\labelenumi{(\alph{enumi})}
\setcounter{enumi}{1}
\item
  再设 A 是 Noether 的；设 q 是属于 \(\operatorname{Ass}(B/\mathfrak{f})\) 的 B 的素理想。证明：若 \(\omega\) 是从 B 到 \(B/q\) 的分式域的典范同态，则 \(\omega(t)\) 在 \(\mathbf{A}/(\mathfrak{q} \cap \mathbf{A})\) 的分式域上是超越的。（若 \(p = q \cap A[t]\)，则存在元素 \(y \in B\)，使 q 是由那些 \(z \in B\)（满足 \(zy \in \mathfrak{f}\)）组成之集；证明 p 是 \(B' = A[t] + By\) 相对于 A{[}t{]} 的导子。然后利用 (a) 以反证法论证。）
\item
  设 A 是 Noether 的。设 n 是 B 的一个极大理想，\(p = A \cap n\)。证明：若 \(A_{p} \neq B_{n}\)，则存在 B 的素理想，它含于 n、异于 n 且包含 p。（设 \(\mathfrak{f}\) 是 B 相对于 A{[}t{]} 的导子。先设 \(n \supset \mathfrak{f}\)；在 B 的含于 n 且包含 \(\mathfrak{f}\) 的素理想所成之集中考虑一个极小元 q；考虑 A 在 B/q 中的典范像，并利用 (b) 与习题 4。若 \(n \not\supset \mathfrak{f}\)，则局部环 B\_n 等于局部环 A{[}t{]}\_\{n\_1\}，其中 n\_1 = n ∩ A{[}t{]}（§ 1 习题 16）。若 t ∈ A\_p，证明 A\_p = B\_n。反之，若 A\_p ≠ B\_n，则由 § 2 习题 14 (c) 推出 pA\_p{[}t{]} 是 A\_p{[}t{]} 的一个非极大素理想，并断定 pB\_n 是 B\_n 的异于 nB\_n 的素理想。）
\end{enumerate}

\begin{enumerate}
\def\labelenumi{\arabic{enumi}.}
\setcounter{enumi}{5}
\tightlist
\item
  若整环 A 是 Noether 的，且对 A 的分式域的扩域中元素的每个有限序列 \((t_{i})_{1\leq i\leq n}\)，\(A[t_{1},\ldots,t_{n}]\) 的整闭包都是有限生成 \(A[t_{1},\ldots,t_{n}]\) -模，则称 A 为整 Noether 整环。域上的每个有限生成整代数都是整 Noether 整环（第 2 节定理 2）。
\end{enumerate}

\begin{enumerate}
\def\labelenumi{(\alph{enumi})}
\item
  整 Noether 整环上的每个有限生成整代数都是整 Noether 整环。不约化为 0 的整 Noether 整环的每个分式环都是整 Noether 整环。
\item
  设 A 是整 Noether 整环，\((t_{i})_{1\leq i\leq n}\) 是 A 的分式域 K 的扩域中元素的一个有限序列，B 是 \(\mathbf{K}(t_{1},\ldots,t_{n})\) 的子环，包含 A 且在 A 上整。则 B 是有限生成 A-模。（利用《代数》第五章 § 5 第 7 节，注意 B 的分式域是 K 的有限代数扩张。）
\end{enumerate}

\begin{enumerate}
\def\labelenumi{\arabic{enumi}.}
\setcounter{enumi}{6}
\item
  设 A 是整 Noether 整环（习题 6），\((t_i)_{1 \leq i \leq n}\) 是 A 的分式域的扩域中元素的一个有限序列，B 是环 \(A[t_1, \ldots, t_n]\)，\(n\) 是 B 的一个极大理想，\(p = A \cap n\)，\(A'\) 是 A 在 B 中的整闭包，\(p' = n \cap A'\)。那么，若局部环 \(A_p'\) 与 \(B_n\) 不同，则存在 B 的一个素理想，含于 \(n\)、异于 \(n\) 且包含 \(p\)（“Zariski 主定理”）。（利用习题 6 (b)，化归为 \(A' = A\) 的情形，再化归为 A 是以 \(p\) 为极大理想的局部环的情形。证明此时 \(B/n\) 是 \(A/p\) 的有限代数扩张（参见第 4 节定理 3）；由此推出，对每个满足 \(A \subset C \subset B\) 的环 C，\(n \cap C\) 都是 C 的极大理想。然后对 \(n\) 用归纳法，并用 \(B_t\) 表示 \(A[t_1, \ldots, t_i]\) 在 B 中的整闭包；按 \(t_{i+1}\) 在 \(B_t\) 的分式域上是超越的还是代数的区分两种情形；在前一情形用习题 4，在后一情形用习题 5 (c)。）
\item
  设 A 是环。为使 A 是 Jacobson 环，必须且只需：对每个极大理想 \(m'\)（多项式环 A{[}X{]} 的），\(m' \cap A\) 都是 A 的极大理想。（注意：若整环 B 有 Jacobson 根 \(R \neq 0\)，且 \(b \in R\)，\(b \neq 0\)，则 B 与 B{[}X{]} 的包含 \(1 + bX\) 的一个极大理想之交不能包含 b。）
\item
  给出一个 Jacobson 整环 A 与包含 A 的 Jacobson 环 B 的例子，使得存在 B 的一个极大理想，它与 A 之交不是极大理想。
\end{enumerate}

¶ 10. 设 k 是域，L 是无限集，A 是多项式环 \(k[X_{\lambda}]_{\lambda \in L}\)。若 \(k_{0}\) 是 k 中所含的素域，用 b 表示 k 在 \(k_{0}\) 上的一个超越基的基数，并记 \(c = \text{Card}(L)\)。

\begin{enumerate}
\def\labelenumi{(\alph{enumi})}
\item
  对 A 的每个理想 a，证明存在 a 的一个生成系 S，使 \(\operatorname{Card}(S) \leqslant c\)。（对 L 的有限子集 J，考虑 a 与子环 \(k[X_{\lambda}]_{\lambda \in J}\) 的交。）
\item
  设 \(c < b\)。证明：若 \(m\) 是 \(A\) 的一个极大理想，则 \(A / m\) 是 \(k\) 的代数扩张。（设 \(K\) 是 \(k\) 的在 \(k_0\) 上由 \(m\) 的一个基数 \(\leqslant c\) 的生成系中诸多项式的系数生成的子域；设 \(m_0 = m \cap K[X_\lambda]_{\lambda \in L}\)，于是 \(m = m_0 \otimes_k k\)，且
\end{enumerate}

\[
\mathrm{A} / \mathfrak {m} = (\mathrm{K} [ \mathrm{X} _ {\lambda} ] _ {\lambda \in \mathrm{L}} / \mathfrak {m} _ {0}) \otimes_ {\mathrm{K}} k.
\]

最后注意，存在从 \(K[X_{\lambda}]_{\lambda \in L}/m_{0}\) 的分式域到 k 的一个代数闭包上的 K-同构。）由此推出，此时 A 是 Jacobson 环（利用习题 8）。

\begin{enumerate}
\def\labelenumi{(\alph{enumi})}
\setcounter{enumi}{2}
\tightlist
\item
  设 \(b < c\)；则 \(\operatorname{Card}(k) \leqslant c\)，从而存在双射 \(\lambda \mapsto \xi_{\lambda}\)，把一个子集 \(L'\)（\(L\) 的）映到 \(k\) 上，且 \(L'\) 在 \(L\) 中的余集非空。设 \(\lambda_0 \in L - L'\)；证明存在一个 \(k\) -同态 \(u\)，它把 \(A\) 映到有理函数域 \(k(Y)\) 的一个子域上，使得 \(u(X_{\lambda_0}) = Y\)，
\end{enumerate}

\[
u (\mathbf {X} _ {\lambda}) = 1 / (\mathbf {Y} - \xi_ {\lambda})
\]

（\(\lambda \in \mathbf{L}'\)），且 \(u(\mathbf{X}_{\lambda}) = 0\)（对 \(\lambda \in \mathbf{L} - (\mathbf{L}' \cup \{\lambda_0\})\)）；若 \(\mathfrak{m}\) 是 \(u\) 的核，则 \(\mathbf{A} / \mathfrak{m}\) 因此不是 \(k\) 的代数扩张。

\begin{enumerate}
\def\labelenumi{(\alph{enumi})}
\setcounter{enumi}{3}
\tightlist
\item
  在与 (c) 相同的假设下，证明 A 不是 Jacobson 环。（注意存在不是 Jacobson 环的 k-代数 B（例如局部环），使 \(\operatorname{Card}(B) = b\)。）
\end{enumerate}


\chapter{赋值}

除另有说明外，本章所考虑的环都假设是交换的且具有单位元素。所有环同态都假设把单位元素映为单位元素。环 A 的每个子环都假设包含 A 的单位元素。

若 A 是局部环，用 \(\mathfrak{m}(\mathbf{A})\) 表示它的极大理想，它的剩余域 \(\mathbf{A}/\mathfrak{m}(\mathbf{A})\) 用 \(\kappa(\mathbf{A})\) 表示， \(\mathbf{U}(\mathbf{A})\) 表示 A 的可逆元素所成的乘法群；于是 \(\mathbf{U}(\mathbf{A}) = \mathbf{A} - \mathfrak{m}(\mathbf{A})\)。

\section{赋值环}

\subsection{局部环之间的支配关系}

定义 1. 设 A、B 是两个局部环。若 A 是 B 的子环且 \(\mathfrak{m}(\mathbf{A}) = \mathbf{A} \cap \mathfrak{m}(\mathbf{B})\)，则称 B 支配 A。

命题 1. 设 A、B 是局部环，且 A 是 B 的子环。下列条件等价：

\begin{enumerate}
\def\labelenumi{(\alph{enumi})}
\item
  \(\mathfrak{m}(\mathbf{A}) \subset \mathfrak{m}(\mathbf{B});\)
\item
  B 支配 A；
\item
  理想 \(\mathbf{B}\mathfrak{m}(\mathbf{A})\)（由 \(\mathfrak{m}(\mathbf{A})\) 在 \(\mathbf{B}\) 中生成）不包含 1。
\end{enumerate}

若 \(\mathfrak{m}(\mathbf{A}) \subset \mathfrak{m}(\mathbf{B})\)，则 \(\mathfrak{m}(\mathbf{B}) \cap \mathbf{A}\) 是 A 的一个不包含 1 且包含极大理想 \(\mathfrak{m}(\mathbf{A})\) 的理想；因此它等于该理想，即 (a) 蕴涵 (b)。若 B 支配 A，则理想 \(\mathbf{B}\mathfrak{m}(\mathbf{A})\) 含于 \(\mathfrak{m}(\mathbf{B})\)，从而不包含 1；于是 (b) 蕴涵 (c)。若 (c) 成立，则 \(\mathbf{B}\mathfrak{m}(\mathbf{A})\) 含于 B 的唯一极大理想 \(\mathfrak{m}(\mathbf{B})\)，由此得 (a)。

注意，若 K 是环，则关系“B 支配 A”是 K 的局部子环所成之集上的一个序关系。

设 A、B 是两个局部环，且 B 支配 A。典范单射 A → B 通过取商定义了从域 κ(A) 到 κ(B) 的一个子域上的同构；借助于这一同构，我们可以把 κ(A) 等同于 κ(B) 的一个子域。

\textbf{例}\par

\begin{enumerate}
\def\labelenumi{(\arabic{enumi})}
\item
  设 A 是 Noether 局部环，\(\hat{\mathbf{A}}\) 是它的完备化；则局部环 \(\hat{\mathbf{A}}\) 支配 A（第三章 § 3 第 5 节命题 9）。
\item
  设 B 是整环，A 是 B 的子环，\(p'\) 是 B 的素理想，且 \(p = A \cap p'\)。则 \(pA_{p} \subset p'B_{p'}\)，从而 \(B_{p'}\) 支配 \(A_{p}\)。
\end{enumerate}

\subsection{赋值环}

定理 1. 设 K 是域，V 是 K 的子环。下列条件等价：

\begin{enumerate}
\def\labelenumi{(\alph{enumi})}
\item
  V 是 K 的局部子环所成之集（该集以 A 与 B 之间的关系“B 支配 A”排序）的极大元。
\item
  存在代数闭域 L 与从 V 到 L 的同态 h，它在从 K 的子环到 L 的同态所成之集（该集以 f 与 g 之间的关系“g 是 f 的扩张”排序）中是极大的。
\item
  若 \(x \in \mathbf{K} - \mathbf{V}\)，则 \(x^{-1} \in \mathbf{V}\)。
\item
  \(\mathbf{V}\) 的分式域是 \(\mathbf{K}\)，且 \(\mathbf{V}\) 的主理想所成之集按包含关系是全序的。
\item
  V 的分式域是 K，且 V 的理想所成之集按包含关系是全序的。
\end{enumerate}

我们通过证明下列蕴涵关系来证明本定理

\[
(a) \Rightarrow (b) \Rightarrow (c) \Rightarrow (d) \Rightarrow (e) \Rightarrow (a).
\]

设 (a) 成立。则 V 是局部环。设 L 是剩余域 \(\kappa(V)\) 的一个代数闭包，h 是从 V 到 L 的典范同态。设 \(V'\) 是 K 的包含 V 的子环，\(h'\) 是从 \(V'\) 到 L 的延拓 h 的同态。若 \(p'\) 是 \(h'\) 的核，则 \(p' \cap V = m(V)\)；于是（第 1 节例 2）\(V_{p'}'\) 支配 V，这蕴涵 \(V_{p'}' = V\) 及 \(V' = V\)。故 (b) 成立。

设 (b) 成立。设 L 是代数闭域，h 是从 V 到 L 的同态；设 h 在从 K 的子环到 L 的同态所成之集中是极大的；设 p 是 h 的核。由于 \(h(V - p)\) 的元素都在 L 中可逆，h 可以延拓为从 \(V_{p}\) 到 L 的同态（第二章 §2 第 1 节命题 1）；于是 \(V = V_{p}\)，这表明 V 是局部环且 p 是它的极大理想。设 x 是 K 的非零元素；须证元素 x、\(x^{-1}\) 中至少有一个属于 V，也就是（依据 h 的极大性）h 可以延拓到 V{[}x{]} 或 \(V[x^{-1}]\)。若 x 在 V 上整，这由第五章 §2 第 1 节定理 1 的推论 4 得出。若 x 不在 V 上整，我们将使用下面的引理：

引理 1. 设 A 是环 B 的子环，x 是 B 的不在 A 上整的元素；则分式环 \(B_{x}\)（第二章 § 5 第 1 节）不约化为 0，且在子环 A{[}1/x{]}（\(B_{x}\) 的）中存在包含 1/x 的极大理想；此外，若 M 是这些极大理想中的任何一个，则 M 在 A 中的原像是极大理想。

由于 \(x\) 不在 \(A\) 上整，它不是幂零的，因此 \(B_x \neq 0\)；此外，\(x / 1 \notin A[1 / x]\)，否则将有一个形如

\[
x / 1 = a _ {0} / 1 + a _ {1} / x + \dots + a _ {n} / x ^ {n}
\]

的关系（其中 \(n \geqslant 0\)，且 \(a_{i} \in \mathbf{A}\)（对 \(0 \leqslant i \leqslant n\)）），它等价于

\[
x ^ {n + h} - a _ {0} x ^ {n + h - 1} - a _ {1} x ^ {n + h - 2} - \dots - a _ {n} x ^ {h - 1} = 0
\]

（其中 \(h \geqslant 1\) 取适当的值）；但这样一个关系将蕴涵 \(x\) 在 \(A\) 上整，与假设相反。于是，\(A[1/x]\) 的包含 \(1/x\) 的极大理想的存在性由 \(1/x\) 在 \(A[1/x]\) 中不可逆这一事实得出（《代数》第一章 § 8 第 7 节定理 2）。

于是设 \(\mathfrak{M}\) 是 A{[}1/x{]} 的包含 1/x 的一个极大理想；设 \(\phi: \mathrm{A} \to \mathrm{A}[1/x]\)，\(p: \mathrm{A}[1/x] \to \mathrm{A}[1/x]/\mathfrak{M}\) 是典范同态；则

\[
p (\mathrm{A} [ 1 / x ]) = p (\phi (\mathrm{A})) [ p (1 / x) ] = p (\phi (\mathrm{A}))
\]

因为 \(p(1/x) = 0\)；这证明 \(p(\phi(A))\) 是域，从而原像 \(\overline{\phi}^{-1}(M)\) 是 A 的极大理想。

我们以 A = V、B = K 应用此引理；于是存在 \(V[x^{-1}]\) 的包含 \(x^{-1}\) 的极大理想 M，且 \(M \cap V\) 是 V 的极大理想；由于 V 是局部的，有 \(M \cap V = p\)；用 f 表示从 \(V[x^{-1}]\) 到 \(V[x^{-1}] / M\) 上的典范同态，则 \(f(x^{-1}) = 0\)，从而 \(V / p = f(V) = f(V[x^{-1}])\)；由于 h 通过取商定义了从 V/p 到 L 的单同态 \(\bar{h}\)，\(\bar{h} \circ f\) 是从 \(V[x^{-1}]\) 到 L 的延拓 h 的同态。故 (c) 成立。

现在设 (c) 成立。显然 K 是 V 的分式域。设 a、b 是 V 的元素，且 \(Va \notin Vb\)；我们证明 \(Vb \subset Va\)。当 b = 0 时这是真的；否则关系 \(a \notin Vb\) 蕴涵 \(b^{-1}a \notin V\)，从而由 (c) 有 \(a^{-1}b \in V\)，因此 \(Vb \subset Va\)。故 (d) 成立。

设 (d) 成立。设 \(\mathfrak{a}\)、\(\mathfrak{b}\) 是 \(\mathbf{V}\) 的理想，且 \(\mathfrak{a} \notin \mathfrak{b}\)。存在 \(a \in \mathfrak{a}\) 使 \(a \notin \mathfrak{b}\)。对所有 \(b \in \mathfrak{b}\)，有 \(a \notin \mathbf{V}b\)，从而 \(Va \notin Vb\)，因此 \(Vb \subset Va \subset \mathfrak{a}\)（由 (c)），即 \(b \in \mathfrak{a}\)。于是 \(\mathfrak{b} \subset \mathfrak{a}\)，这表明条件 (e) 满足。

最后设 (e) 成立。由于 V 有极大理想，它只有一个，从而是局部环。设 \(V'\) 是 K 的支配 V 的局部子环，x 是 \(V'\) 的非零元素；记 \(x = ab^{-1}\)，其中 \(a \in V\)、\(b \in V\)。理想 Va、Vb 之一含于另一个。若 Va ⊂ Vb，则 \(x \in V\)。若 Vb ⊂ Va，则 \(x^{-1} \in V\)；由于理想 \(V'm(V)\) 不包含 1（第 1 节命题 1），\(x^{-1} \notin m(V)\)，于是仍得 \(x \in V\)，因为 \(V\) 是局部的。因此 \(V'\) 的每个元素都属于 \(V\)；我们断定 (a) 成立。

定义 2. 在定理 1 的记号下，若等价条件 (a)、(b)、(c)、(d)、(e) 成立，则称 V 是域 K 的赋值环。若一个环是整环且是它的分式域的赋值环，则称之为赋值环。

定理 2. 设 K 是域，h 是从 K 的子环 A 到代数闭域 L 的同态。则存在 K 的赋值环 V 与一个同态

\(h^{\prime}\) 从 \(\mathbf{V}\) 到 \(\mathbf{L}\)，使得 \(\mathbf{V}\) 包含 \(\mathbf{A}\)，\(h^{\prime}\) 延拓 \(h\)，且 \(h^{\prime}(0) = \mathfrak{m}(\mathbf{V})\)。

设 \(\mathfrak{S}\) 是从 K 的子环到 L 的同态所成之集，以扩张关系排序。这个集是归纳的；因为若 \((h_{\alpha})_{\alpha \in I}\) 是 \(\mathfrak{S}\) 的元素的一个非空全序族，\(B_{\alpha}\) 是 \(h_{\alpha}\) 的定义环，则诸 \(B_{\alpha}\) 构成 K 的子环的一个全序族，它们的并 B 因而是 K 的子环；于是存在唯一的从 B 到 L 的映射 \(\bar{h}\) 延拓诸 \(h_{\alpha}\)（《集合论》第二章 § 4 第 6 节命题 7），并且立即可见 \(\bar{h}\) 是从 B 到 L 的同态。Zorn 引理于是表明存在一个极大元 \(h'\)，它属于 \(\mathfrak{S}\) 且延拓 \(h\)。\(h'\) 的定义环 V 是 K 的赋值环（定理 1）；若 \(p\) 是 \(h'\) 的核，则 \(h'\) 可以延拓为从 \(V_p\) 到 L 的同态（第二章 § 2 第 1 节命题 1），于是 \(V_p = V\) 且 \(p = m(V)\)。

推论. 域 K 的每个局部子环 A 至少被 K 的一个赋值环所支配。

把定理 2 应用于典范同态 \(h\)，它从 \(\mathbf{A}\) 映到一个代数闭包 \(\mathbf{L}\)（\(\mathbf{A} / \mathfrak{m}(\mathbf{A})\) 的）中。

\subsection{整元素的刻画}

定理 3. 设 A 是域 K 的子环。A 在 K 中的整闭包 A' 是 K 的包含 A 的诸赋值环之交；若 A 是局部的，则 A' 是 K 的支配 A 的诸赋值环之交。

设 \(x\) 是 \(A'\) 的元素，\(V\) 是 \(K\) 的包含 \(A\) 的赋值环；由于 \(x\) 在 \(V\) 上整，存在素理想 \(p'\)（\(V[x]\) 的）使 \(p' \cap V = m(V)\)（第五章 §2 第 1 节定理 1）；显然局部环 \((V[x])_{p'}\) 支配 \(V\)，从而等于它；于是 \(x \in V\)。反之，设 \(y\) 是 \(K\) 的不在 \(A\) 上整的元素；则存在极大理想 \(M\)（\(A[y^{-1}]\) 的），它包含 \(y^{-1}\)（第 2 节引理 1）；还存在赋值环 \(V\)（\(K\) 的），它支配 \((A[y^{-1}])_{m}\)（第 2 节定理 2 的推论）；由于 \(y^{-1} \in m(V)\)，有 \(y \notin V\)。此外 \(M \cap A\) 是 \(A\) 的极大理想（第 2 节引理 1）；因此，若 \(A\) 是局部的，则 \(M \cap A = m(A)\)，从而 \(V\) 支配 \(A\)。

推论 1. 每个赋值环都是整闭的。

推论 2. 为使一个整环是整闭的，必须且只需它是它的分式域的一族赋值环之交。

在 Noether 环的情形，推论 2 可以更精确化（第七章 § 1 第 3 节定理 1 的推论）。

推论 3. 设 K 是域，K' 是 K 的扩张，A 是 K 的赋值环。A 在 K' 中的整闭包是 K' 的满足 V' ∩ K = A 的诸赋值环 V' 之交。

定理 1 (c) 表明，若 \(V'\) 是 \(K'\) 的赋值环，则 \(V' \cap K\) 是 K 的赋值环，且 \(V'\) 支配 \(V' \cap K\)。为使 \(V'\) 支配 A，必须且只需 \(V' \cap K\) 支配 A，从而等于 A。

\subsection{赋值环的例子}

\begin{enumerate}
\def\labelenumi{(\arabic{enumi})}
\item
  每个域都是赋值环。
\item
  若 \(\mathbf{V}'\) 是域 \(\mathbf{K}'\) 的赋值环，\(\mathbf{K}\) 是 \(\mathbf{K}', \mathbf{V}' \cap \mathbf{K}\) 的子域，则由第 2 节定理 1 (c)，后者是 \(\mathbf{K}\) 的赋值环。
\item
  下面的命题给出赋值环的大量例子：
\end{enumerate}

命题 2. 设 A 是局部环，其极大理想是主理想 Ap。若 \(\bigcap_{n=1}^{\infty} A p^n = (0)\)（例如当 A 是 Noether 时，参见第三章 §3 第 2 节

命题 5 的推论），则 \(\mathbf{A}\) 的理想只有 (0) 与诸 \(\mathbf{A}p^n\)；于是或者 \(p\) 是幂零的，或者 \(\mathbf{A}\) 是赋值环。

设 A 被诸 \(\mathbf{A}p^n\) 滤过，用 \(v\) 表示相应的阶函数（第三章 § 2 第 2 节）。由于

\[
\bigcap_ {n = 1} ^ {\infty} A p ^ {n} = (0),
\]

关系 \(v(x) = +\infty\) 蕴涵 \(x = 0\)。设 \(\mathfrak{a}\) 是一个理想 \(\neq (0)\)，属于 \(\mathbf{A}\)，\(a\) 是 \(\mathfrak{a}\) 中使 \(v\) 取到最小值的一个元素；记 \(v(a) = s\)（\(s \neq +\infty\)）。则 \(\mathfrak{a} \subset \mathrm{Ap}^s\)。特别地，存在 \(u \in \mathbf{A}\) 使 \(a = up^s\)；由于 \(a \notin \mathrm{Ap}^{s+1}\)，有 \(u \notin \mathrm{Ap}\)；于是 \(u\) 可逆且 \(p^s \in \mathbf{A}a \subset \mathfrak{a}\)。由此得出 \(\mathfrak{a} = \mathrm{Ap}^s\)，即得我们的第一个断言。同样可见，每个元素 \(a \neq 0\)（\(\mathbf{A}\) 的）都可写成 \(a = up^{v(a)}\) 的形式，其中 \(u\) 可逆。若 \(a' = u'p^{v(a')}\)（\(u'\) 可逆）是 \(\mathbf{A}\) 的另一个非零元素，则 \(aa' = uu'p^{v(a)+v(a')}\)；于是，若 \(p\) 不是幂零的，则 \(aa' \neq 0\)，从而 \(\mathbf{A}\) 是整环。这时，由于 \(\mathbf{A}\) 的理想所成之集按包含关系是全序的，我们断定 \(\mathbf{A}\) 是赋值环（定理 1 (e)）。

例如，若 \(p\) 是素数，则局部环 \(\mathbf{Z}_{(p)}\) 是赋值环。设 \(\mathbf{B} = \mathbf{K}[\mathbf{X}_1, \ldots, \mathbf{X}_n]\) 是 \(n\) 个未定元在域 \(\mathbf{K}\) 上的多项式环；理想 \(\mathbf{BX}_1\) 是素的，因为 \(\mathbf{B} / \mathbf{BX}_1\) 同构于 \(\mathbf{K}[\mathbf{X}_2, \ldots, \mathbf{X}_n]\)；于是 \(\mathbf{B}_{\mathbf{BX}_1}\) 是赋值环；它由有理函数 \(\mathbf{PQ}^{-1}\) 组成，其中 \(\mathbf{P}\)、\(\mathbf{Q}\) 是多项式，且 \(\mathbf{Q}(0, \mathbf{X}_2, \ldots, \mathbf{X}_n) \neq 0\)。

* 更一般地，我们将看到，若 F 是

\[
\mathrm{B} = \mathrm{K} [ \mathrm{X} _ {1}, \dots , \mathrm{X} _ {n} ],
\]

的一个极端元，则 \(B_{BF}\) 是赋值环（参见第七章 § 3 第 5 节）。*

域 K 上一个未定元的形式幂级数环 K{[}{[}T{]}{]} 是 Noether 局部整环，其极大理想是主的；因此它是赋值环。另一方面，两个未定元的形式幂级数环 K{[}{[}T₁, T₂{]}{]} 虽是 Noether 局部整环，却不是赋值环，因为元素 T₁、T₂ 中没有一个是另一个的倍数。

命题 3. 设 A 是主理想整环，K 是它的分式域。K 的包含 A 且异于 K 的赋值环恰是形如 \(A_{Ap}\) 的环，其中 p 是 A 的极端元。

显然 \(A_{Ap}\)（p 为极端元）是包含 A 且异于 K 的赋值环（命题 2）。反之，设 V 是异于 K 且包含 A 的赋值环。由于 \(V \neq K\)，\(m(V)\) 包含一个元素 \(x \neq 0\)；记 x = a/b，其中 \(a \in A\)、\(b \in A\)，可见 \(A \cap m(V)\) 包含非零元素 a。由于 \(A \cap m(V)\) 是素的，它形如 Ap，其中 p 在 A 中是极端的。于是 \(A_{Ap} \subset V\)，\(pA_{Ap} \subset m(V)\)，从而 V 支配 \(A_{Ap}\)；由于 \(A_{Ap}\) 是 K 的赋值环（命题 2），有 \(V = A_{Ap}\)。

推论 1. 域 \(\mathbf{Q}\) 的每个异于 \(\mathbf{Q}\) 的赋值环都形如 \(\mathbf{Z}_{(p)}\)，其中 \(p\) 是素数。

Q 的每个子环都包含 Z。

推论 2. 设 K 是域，K(X) 是 K 上一个未定元的有理函数域，V 是 K(X) 的包含 K 且异于 K(X) 的赋值环。若 X ∈ V，则存在不可约多项式 P ∈ K{[}X{]}，使 \(V = (K{[}X{]})_{(P)}\)；否则 V 是 K{[}X⁻¹{]} 在素理想 X⁻¹K{[}X⁻¹{]} 处的局部环（换言之，V 由分式 A/B 组成，其中 A ∈ K{[}X{]}、B ∈ K{[}X{]}，且 deg(A) ≤ deg(B)）。

若 \(\mathbf{X} \in \mathbf{V}\)，则 \(\mathbf{K}[\mathbf{X}] \subset \mathbf{V}\)，所述断言由命题 3 得出。若 \(\mathbf{X} \notin \mathbf{V}\)，则 \(\mathbf{X}^{-1} \in \mathbf{V}\)，从而 \(\mathbf{K}[\mathbf{X}^{-1}] \subset \mathbf{V}\)；于是 \(\mathbf{V}\) 是 \(\mathbf{K}[\mathbf{X}^{-1}]\) 在某个素理想 \(p\) 处的局部环（命题 3），且这个理想包含 \(\mathbf{X}^{-1}\)，因为 \(\mathbf{X}^{-1}\) 在 V 中不可逆；于是 \(p = X^{-1}K[X^{-1}]\)，因为后一理想是极大的。最后考虑有理函数 \(A(X)/B(X)\)，其中 A、B 是次数分别为 a 与 b 的多项式；则 \(A(X) = X^{a}A'(X^{-1})\)，\(B(X) = X^{b}B'(X^{-1})\)，其中 \(A'\)、\(B'\) 是满足 \(A'(0) \neq 0\) 与 \(B'(0) \neq 0\) 的多项式；于是，为使 \(A(X)/B(X)\) 属于 \(K[X^{-1}]\) 在 \(X^{-1}K[X^{-1}]\) 处的局部环，必须且只需 \(a \leqslant b\)。

\section{位}

\subsection{非处处有定义的合成律的态射概念}

定义 1. 设 E、E' 是两个集合，各带一个记作 \((x,y)\mapsto x*y\) 的内合成律，它不必处处有定义。映射 \(f:E\to E'\) 称为态射，如果对 E 中所有使 \(f(x)*f(y)\) 有定义的 x、y，合成 \(x*y\) 也有定义，且：

\[
f (x * y) = f (x) * f (y).\tag{1}
\]

更简短地说，公式 (1) 必须在右端有意义的每一次都成立。

态射的概念不同于表示的概念（《代数》

第一章 § 1 第 1 节），在那里要求只要左端有意义，等式 (1) 就成立。当然，对于处处有定义的合成律，这两个概念是一致的。

定义 2. 设 E、E' 是两个集合，各带一族内合成律 \((x,y)\mapsto x*\alpha y,\alpha\in\mathbf{I}\)。映射 \(f:E\to E'\) 称为态射，如果它对每个合成律 \((x,y)\mapsto x*\alpha y\) 都是态射。

与表示一样，态射也满足《集合论》第四章 § 2 的公理 \((\mathrm{MO}_{\mathrm{I}})\)、\((\mathrm{MO}_{\mathrm{II}})\)、\((\mathrm{MO}_{\mathrm{III}})\)。若 \(f: E \to E'\) 是态射，则 \(f(E)\) 是 \(R'\) 的稳定子集。

\subsection{位}

若 K 是域，回忆 \(\tilde{K}\) 表示 K 与一个记作 \(\infty\) 的元素之和（的集）（《代数》第二章 § 9 第 9 节）；K 的合成律按如下方式（出处同上）扩张到 \(\tilde{K}\)

\begin{enumerate}
\def\labelenumi{(\arabic{enumi})}
\setcounter{enumi}{1}
\tightlist
\item
\end{enumerate}

\[
a + \infty = \infty \quad \text { for } \quad a \in \mathbf {K}, \quad a \neq \infty ,\tag{2}
\]

\[
\infty . a = a. \infty = \infty \quad \text {   for   } \quad a \in \tilde {\mathbf {K}}, \quad a \neq 0.\tag{3}
\]

于是没有定义的合成只有 \(\infty + \infty\)、\(\infty, 0\) 与 \(0, \infty\)。另一方面，映射 \(x \mapsto -x\) 与 \(x \mapsto x^{-1}\) 类似地扩张到 \(\tilde{\mathbf{K}}\)，通过令 \(-\infty = \infty, 0^{-1} = \infty, \infty^{-1} = 0\)。我们还将记 \(x + (-y) = x - y\)。

集 \(\tilde{\mathbf{K}}\) 称为与 \(\mathbf{K}\) 相伴的射影域，它可等同于射影直线 \(\mathbf{P}_1(\mathbf{K})\)（出处同上）。

定义 3. 设 K、L 是两个域。\(\tilde{\mathbf{K}}\) 到 \(\tilde{\mathbf{L}}\) 的（关于加法与乘法的）每个满足 \(f(1) = 1\) 的态射 f 称为 K 的以 L 为值域的位。

换言之，若 \(x\)、\(y\) 是 \(\tilde{\mathbf{K}}\) 的元素，且 \(f(x) + f(y)\)（相应地，\(f(x)f(y)\)）有定义，则 \(x + y\)（相应地，\(xy\)）有定义，且

\begin{enumerate}
\def\labelenumi{(\arabic{enumi})}
\setcounter{enumi}{3}
\tightlist
\item
\end{enumerate}

\[
f (x + y) = f (x) + f (y)\tag{4}
\]

\[
f (x y) = f (x) f (y).\tag{5}
\]

由于 \(\infty +\infty\) 没有定义，\(f(\infty) + f(\infty)\) 也没有定义，这表明

\[
f (\infty) = \infty .\tag{6}
\]

类似地，由于 \(0.\infty\) 没有定义，\(f(0)f(\infty)\) 也没有定义，由 (6) 这蕴涵

\[
f (0) = 0.\tag{7}
\]

另一方面

\[
f (a ^ {- 1}) = f (a) ^ {- 1} \quad \text {   for   all   } a \in \tilde {\mathbf {K}}.\tag{8}
\]

若 \(f(a)f(a^{-1})\) 有定义，则 \(aa^{-1}\) 有定义，从而等于 1；于是 \(f(a)f(a^{-1}) = f(1) = 1\)，这在此情形证明了 (8)。若 \(f(a)f(a^{-1})\) 没有定义，则或者 \(f(a) = 0\) 且 \(f(a^{-1}) = \infty\)，或者 \(f(a) = \infty\) 且 \(f(a^{-1}) = 0\)，而 (8) 仍成立。

类似地可以证明

\[
f (- a) = - f (a) \quad \text {   for   all   } a \in \tilde {\mathbf {K}}.\tag{9}
\]

由公式 (8) 与 (9) 得出，\(f\) 对合成律 \((x, y) \mapsto x - y\) 与 \((x, y) \mapsto xy^{-1}\) 也是态射。

对 \(x \in \tilde{\mathbf{K}}\)，称 \(f\) 在 \(x\) 处有限，如果 \(f(x) \neq \infty\)；由 (6) 这蕴涵 \(x \in \mathbf{K}\)。

若 \(f \colon \tilde{\mathbf{K}} \to \tilde{\mathbf{L}}\) 是位，则 \(f(\tilde{\mathbf{K}})\) 是 \(\tilde{\mathbf{L}}\) 的一个子集，它对合成律 \((x, y) \mapsto x + y\)、\((x, y) \mapsto x - y\)、\((x, y) \mapsto xy\) 与 \((x, y) \mapsto xy^{-1}\) 稳定，且包含 1。若 \(\mathbf{E}\) 是 \(f(\tilde{\mathbf{K}})\) 的有限元素所成之集，则 \(\mathbf{E}\) 是 \(\mathbf{L}\) 的子域，且 \(f(\tilde{\mathbf{K}}) = \tilde{\mathbf{E}}\)。习惯上称 \(\mathbf{E}\) 为 \(f\) 的值域。

两个位的合成映射是位。

设 F 是从域 K 到域 L 的一个子域上的同构；把 f 扩张到 \(\tilde{K}\)，令 \(f(\infty) = \infty\)。如此得到 K 的一个以 L 为值域的位，称之为平凡位，它常与同构 f 等同。

\subsection{位与赋值环}

命题 1. 设 K 是域，A 是 K 的赋值环，\(\kappa(A)\) 是 A 的剩余域。把从 A 到 \(\kappa(A)\) 上的典范映射扩张为映射 \(h_{A}: \tilde{K} \to (\kappa(A))^{\sim}\)，依据等式 \(h_{A}(x) = \infty\)（当 \(x \notin A\)）。如此定义的映射 \(h_{A}\) 是 K 的一个位，其值域为 \(\kappa(A)\)。

显然 \(h_{\mathbf{A}}(1) = 1\)。

我们证明 \(h_{\mathbf{A}}\) 是关于加法的态射。设 \(x, y\) 是 \(\tilde{\mathbf{K}}\) 的两个元素，使 \(h_{\mathbf{A}}(x) + h_{\mathbf{A}}(y)\) 有定义。这时两个元素 \(x, y\) 之一属于 \(\mathbf{A}\)，从而 \(x + y\) 有定义。若 \(x \in \mathbf{A}\) 且 \(y \in \mathbf{A}\)，显然

\[
h _ {\mathrm{A}} (x) + h _ {\mathrm{A}} (y) = h _ {\mathrm{A}} (x + y)
\]

成立。若 \(x \in \mathbf{A}\) 而 \(y \notin \mathbf{A}\)，则 \(x + y \notin \mathbf{A}\)，上式的两端都等于 \(\infty\)。

最后证明 \(h_{\mathbf{A}}\) 是关于乘法的态射。设 \(x \in \tilde{\mathbf{K}}, y \in \tilde{\mathbf{K}}\) 使 \(h_{\mathbf{A}}(x)h_{\mathbf{A}}(y)\) 有定义。若 \(x \in \mathbf{A}\) 且 \(y \in \mathbf{A}\)，显然 \(xy\) 有定义且 \(h_{\mathbf{A}}(x)h_{\mathbf{A}}(y) = h_{\mathbf{A}}(xy)\)。现在设元素 \(x, y\) 之一（例如 \(y\)）不属于 \(\mathbf{A}\)；由于 \(h_{\mathbf{A}}(y) = \infty\)，有 \(h_{\mathbf{A}}(x) \neq 0\)，即 \(x \notin \mathfrak{m}(\mathbf{A})\)，从而 \(x^{-1} \in \mathbf{A}\)；由此得出 \(xy\) 有定义且 \(xy \notin \mathbf{A}\)，于是

\[
h _ {\mathrm{A}} (x y) = \infty = h _ {\mathrm{A}} (x) h _ {\mathrm{A}} (y).
\]

这就证明了命题 1。

若 \(j\) 是从 \(\kappa(A)\) 到某个域 \(L, j \circ h_{A}: \tilde{K} \to \tilde{L}\) 的一个子域上的同构，则它是 \(K\) 的以 \(L\) 为值域的位。上述过程实际上给出了 \(K\) 上的所有位。确切地说：

命题 2. 设 K、L 是两个域，f 是 K 的以 L 为值域的位。则存在 K 的赋值环 A 与从 \(\kappa(\mathbf{A})\) 到 L 的一个子域上的同构 j，使 \(f = j \circ h_{A}\)；这些条件唯一确定 A 与 j。环 A 是由 \(x \in K\) 中使 \(f(x) \neq \infty\) 的元素组成之集，\(\mathfrak{m}(\mathbf{A})\) 是由 \(x \in K\) 中使 \(f(x) = 0\) 的元素组成之集。

若 \(f = j \circ h_{\mathbf{A}}\)，则条件 \(f(x) \neq \infty\)（相应地，\(f(x) = 0\)）等价于条件 \(h_{\mathbf{A}}(x) \neq \infty\)（相应地，\(h_{\mathbf{A}}(x) = 0\)），从而等价于条件 \(x \in \mathbf{A}\)（相应地，\(x \in \mathfrak{m}(\mathbf{A})\)）。于是 \(\mathbf{A}\) 被唯一确定，且由于 \(h_{\mathbf{A}}\) 是满射，\(j\) 也是唯一的。

现在设 \(f\) 是 \(\mathbf{K}\) 的以 \(\mathbf{L}\) 为值域的任一位；用 \(\mathbf{A}\) 表示由 \(x \in \mathbf{K}\) 中使 \(f(x) \neq \infty\) 的元素组成之集。若 \(x \in \mathbf{A}\) 且 \(y \in \mathbf{A}\)，则合成 \(f(x) - f(y)\) 与 \(f(x)f(y)\) 有定义且 \(\neq \infty\)，这表明 \(x - y \in \mathbf{A}\) 且 \(xy \in \mathbf{A}\)；于是 \(\mathbf{A}\) 是 \(\mathbf{K}\) 的子环。若 \(x \notin \mathbf{A}\)，则 \(f(x) = \infty\)，从而 \(f(x^{-1}) = 0\)，于是 \(x^{-1}\) 属于核 \(m\)，其中 \(f'\) 是把 \(f\) 限制到 \(\mathbf{A}\) 所得的同态。反之，若 \(y \in m\)，则 \(y^{-1} \notin \mathbf{A}\)。这表明 \(\mathbf{A}\) 是 \(\mathbf{K}\) 的赋值环，且 \(m\) 是它的极大理想。设 \(j\) 是从 \(\kappa(\mathbf{A})\) 到 \(L\) 的单同态，它是由 \(f'\) 通过取商导出的。则 \(f(x) = j(h_{\mathrm{A}}(x))\)（对所有 \(x \in A\)），且这个等式对 \(x \notin A\) 仍成立，此时两端都等于 \(\infty\)。

分解 \(f = j \circ h_{A}\) 称为位 f 的典范分解。A 称为 f 的环，\(\mathfrak{m}(A)\) 称为 f 的理想，\(\kappa(A)\) 称为 f 的剩余域。为使两个位 \(f: \tilde{K} \to \tilde{L}\) 与 \(f': \tilde{K} \to \tilde{L}'\) 有相同的环，必须且只需存在从 f 的值域到 \(f'\) 的值域上的同构 s，使 \(f' = s \circ f\)；这时称 f 与 \(f'\) 等价。可以看出，关于赋值环的每个结果都可以翻译成关于位的结果，反之亦然；下面各节我们要做的正是这件事。

\textbf{位的例子}\par

\begin{enumerate}
\def\labelenumi{(\arabic{enumi})}
\item
  设 \(\mathbf{K}\) 是域。\(\mathbf{K}\) 上的恒等映射是平凡位，其环为 \(\mathbf{K}\)，理想为 (0)。
\item
  设 \(k\) 是域。对所有 \(u \in k((\mathrm{T}))^{\sim}\)，记 \(f(u) = \infty\)（当 \(u \notin k[[\mathrm{T}]]\) 时），并定义 \(f(u)\) 为 \(u\) 的常数项（当 \(u \in k[[\mathrm{T}]]\) 时）。则 \(f\) 是 \(k((\mathrm{T}))\) 的一个位，其剩余域为 \(k\)，环为 \(k[[\mathrm{T}]]\)。因为 \(k[[\mathrm{T}]]\) 是 \(k((\mathrm{T}))\) 的赋值环（§ 1 第 4 节例 3），且 \(f\) 在 \(k[[\mathrm{T}]]\) 上的限制等同于从 \(k[[\mathrm{T}]]\) 到它的剩余域上的典范同态。
\item
  设 \(k\) 是域，\(a\) 是 \(k\) 的元素，\(A\) 是由这样的 \(u \in k(X)\) 组成之集：\(a\) 可代入 \(u\)（《代数》第四章 § 3 第 2 节）。若用 \(p\) 表示素理想 \((X - a)\)（\(k[X]\) 的），则 \(A = k[X]_p\)，从而 \(A\) 是 \(k(X)\) 的赋值环（§ 1 第 4 节命题 2）。对所有 \(u \in k(X)^{\sim}\)，记 \(f(u) = \infty\)（当 \(u \notin A\) 时），记 \(f(u) = u(a)\)（当 \(u \in A\) 时）。则 \(f\) 是 \(k(X)\) 的一个位，其剩余域为 \(k\)，环为 \(A\)；因为 \(f\) 在 \(A\) 上的限制是 \(A\) 到 \(k\) 上的同态（《代数》第四章 § 3 命题 2），其核为 \(pA = m(A)\)。称元素 \(f(u) \in \tilde{k}\) 是由令 \(X = a\)（在 \(u\) 中）而得到的。
\end{enumerate}

* (4) 设 S 是一维连通复解析簇，K 是 S 上的亚纯函数域。对所有 \(z_0 \in \mathbf{S}\)，映射 \(f \mapsto f(z_0)\)（从 K 到 \(\tilde{\mathbf{C}}\)）是 K 的一个位，其环是 \(f \in \mathbf{K}\) 中在 \(z_0\) 处全纯的那些元素所成之集，其理想是 \(f \in \mathbf{K}\) 中在 \(z_0\) 处为零的那些元素所成之集。正是这个例子以及其他类似的例子，成为术语“位”的来源。*

\subsection{位的扩张}

命题 3. 设 K 是域，S 是 K 的子环，f 是从 S 到代数闭域 L 的同态。则存在 K 的以 L 为值域的延拓 f 的位。

把命题 1 考虑在内，这就是 §1 第 2 节定理 2 的一个翻译。

命题 4. 设 K 是域，f 是 K 的以域 L 为值域的位，K' 是 K 的扩张。则存在 L 的扩张 L' 与 K' 的以 L' 为值域的延拓 f 的位 f'。若 \(x_{1}, \ldots, x_{n}\) 是 K' 的在 K 上代数无关的元素，\(a_1, \ldots, a_n\) 是 L 的任意元素，则可如此选取 \(f'\)，使 \(f'(x_i) = a_i\)（对所有 \(1 \leqslant i \leqslant n\)）。

设 V 是 f 的环，g 是 f 在 V 上的限制，\(g'\) 是 g 的到 \(V[x_{1},\ldots,x_{n}]\) 上的满足 \(g'(x_{i}) = a_{i}\)（对所有 \(1 \leqslant i \leqslant n\)）的扩张。只需取 \(L'\) 为 L 的一个代数闭包，并对 \(g'\) 与 \(L'\) 应用命题 3：我们得到一个位 \(f'\colon \tilde{K}' \to \tilde{L}'\)，它延拓 \(g'\)；若 \(x \in \tilde{K} - V\)，则 \(x^{-1} \in m(V)\)，从而 \(f(x^{-1}) = g(x^{-1}) = 0\)，且 \(f'(x) = \infty = f(x)\)；于是 \(f'\) 延拓 f。

\subsection{用位刻画整元素}

命题 5. 设 K 是域，S 是 K 的子环，h 是从 S 到某个域的同态，p 是 h 的核。为使 K 的元素 x 在局部环 \(S_{p}\) 上整，必须且只需 K 的每个延拓 h 的位都在 x 处有限。

若 \(f\) 是 \(K\) 的延拓 \(h, f\) 在 \(S_p\) 上有限、在 \(pS_p\) 上为零，从而位 \(f\) 的环支配 \(S_p\)。反之，若 \(V\) 是 \(K\) 的支配 \(S_p\) 的赋值环，则 \(V\) 是某个位 \(f\) 的环，它在 \(S\) 上的限制是与 \(h\) 有相同核的同态；把 \(f\) 换成一个等价的位，即可看出 \(V\) 是 \(K\) 的某个延拓 \(h\) 的位的环。说 \(K\) 的每个延拓 \(h\) 的位都在 \(x\) 处有限，等价于说 \(x\) 属于 \(K\) 的支配 \(S_p\) 的所有赋值环。于是命题由 § 1 第 3 节定理 3 得出。

命题 6. 设 K 是域，S 是 K 的子环。为使元素 \(x \in \mathbf{K}\) 在 S 上整，必须且只需 K 的每个在 S 上有限的位都在 \(x\) 处有限。

这同样是 §1 第 3 节定理 3 的推论。

\section{赋值}

\subsection{环上的赋值}

设 \(\Gamma\) 是以加法记写的全序交换群。在本章的其余部分，对这样一个群，我们须考虑把一个元素添加到 \(\Gamma\) 上所得到的集，该元素记作 \(+\infty\)；把这个集记作 \(\Gamma_{\infty}\)，并赋予它：(1) 一个使 \(+\infty\) 为最大元素的全序，换言之，使得 \(\alpha < +\infty\)（对所有 \(\alpha \in \Gamma\)）；(2) 一个交换幺半群结构，其法则在 \(\Gamma\) 上诱导出给定的群法则，且由等式

\[
(+ \infty) + (+ \infty) = + \infty , \quad \alpha + (+ \infty) = + \infty
\]

（对所有 \(\alpha \in \Gamma\)）定义；立即可验证这个法则是结合的、交换的，且关系 \(\alpha \leqslant \beta\)（在 \(\Gamma_{\infty}\) 中）蕴涵 \(\alpha + \gamma \leqslant \beta + \gamma\)（对所有 \(\gamma \in \Gamma_{\infty}\)）。

定义 1. 设 C 是一个（不必交换的）环，\(\Gamma\) 是以加法记写的全序交换群。C 的以 \(\Gamma\) 为值群的任一赋值是指满足下列条件的任一映射 \(v: \mathbf{C} \to \Gamma_{\infty}\)：

\[
(\mathrm{VL}_{\mathrm{I}}) \quad v (x y) = v (x) + v (y) for x \in \mathbf {C}, y \in \mathbf {C}.
\]

\[
(\mathrm{VL}_{\mathrm{II}}) \quad v (x + y) \geqslant \inf (v (x), v (y)) for x \in \mathbf {C}, y \in \mathbf {C}.
\]

\[
(\mathrm{VL} _ {\mathrm{III}}) \quad v (1) = 0 and v (0) = + \infty .
\]

若 C 除 0 外没有零因子，则唯一的映射 \(v_0\)（从 C 到 \(\Gamma_\infty\)）——满足 \(v_0(x) = 0\)（对 \(x \neq 0\)）且 \(v_0(0) = +\infty\)——是一个赋值，称为 C 上的非真赋值。若 \(z \in \mathbf{C}\) 满足 \(z^n = 1\)（对某个整数 \(n \geqslant 1\)），则由 \((\mathrm{VL}_{\mathrm{I}})\)，\(nv(z) = v(z^n) = 0\)，从而 \(v(z) = 0\)（对 C 上的每个赋值 \(v\)），因为 \(\Gamma\) 是全序群。特别地 \(v(-1) = 0\)，于是 \(v(-x) = v(x)\)（对所有 \(x \in \mathbf{C}\)）。此外，由 \((\mathrm{VL}_{\mathrm{I}})\) 得出 \(v(xy) = v(yx)\)（对 C 中所有 \(x, y\)）。若 \(x\) 在 C 中可逆，则 \(v(x^{-1}) = -v(x)\)。

命题 1. 设 \(v\) 是 \(a\)（不必交换的）环 \(\mathbf{C}\) 上的赋值。对任意元素 \(x_{i} \in \mathbf{C}\)（\(1 \leqslant i \leqslant n\)），

\[
v \left(\sum_ {i = 1} ^ {n} x _ {i}\right) \geqslant \inf _ {1 \leqslant i \leqslant n} v (x _ {i})\tag{1}
\]

此外，若存在唯一的指标 \(k\) 使 \(v(x_k) = \inf_{1 \leqslant i \leqslant n} v(x_i)\)，则 (1) 的两端相等。特别地，若 \(v(x) \neq v(y)\)，则 \(v(x + y) = \inf(v(x), v(y))\)。

关系 (1) 由公理 \((\mathrm{VL}_{\mathrm{II}})\) 对 \(n\) 用归纳法得出。若存在唯一的指标 \(k\) 使 \(v(x_k) = \inf_{1 \leqslant i \leqslant n} v(x_i)\)，则记 \(y = \sum_{i \neq k} x_i\)、\(z = \sum_{i=1}^{n} x_i\)，由 (1) 有 \(v(y) > v(x_k)\) 且 \(v(z) \geqslant v(x_k)\)；若 \(v(z) > v(x_k)\)，则关系 \(x_k = z - y\) 将给出 \(v(x_k) \geqslant \inf(v(z), v(y)) > v(x_k)\)，这是荒谬的；于是 \(v(z) = v(x_k)\)，这就证明了第二个断言。

推论. 若 C 的元素的一个有限序列 \((x_{i})_{1\leqslant i\leqslant n}\)（其中 \(n\geqslant 2\)）满足 \(\sum_{i = 1}^{n}x_{i} = 0\)，则至少存在两个不同的指标 \(j,k\)，使

\[
v (x _ {j}) = v (x _ {k}) = \inf _ {1 \leqslant i \leqslant n} v (x _ {i}).
\]

若只有唯一的指标 \(k\) 使 \(v(x_k) = \inf_{1 \leqslant i \leqslant n} v(x_i)\)，则命题 1 将表明 \(v(x_k) = v(0) = +\infty\)，从而 \(v(x_i) = +\infty\)（对所有 \(i\)），这与关系 \(n \geqslant 2\) 及对 \(k\) 所作的假设相反。

注

\begin{enumerate}
\def\labelenumi{(\arabic{enumi})}
\item
  若 \(v: \mathbf{C} \to \Gamma_{\infty}\) 是 \(\mathbf{C}\) 上的赋值，\(u: \mathbf{B} \to \mathbf{C}\) 是从环 \(\mathbf{B}\) 到 \(\mathbf{C}\) 的同态，则立即可见合成映射 \(\mathbf{B} \xrightarrow{u} \mathbf{C} \xrightarrow{v} \Gamma_{\infty}\) 是 \(\mathbf{B}\) 的以 \(\Gamma\) 为值群的赋值。
\item
  条件 \((\mathrm{VL}_{\mathrm{I}})\) 与 \((\mathrm{VL}_{\mathrm{II}})\) 立即表明，集 \(\overline{v}^{-1}(+\infty)\) 是 C 中的一个双边理想 p，且由 \((\mathrm{VL}_{\mathrm{III}})\)，它异于 C；此外，若 x、y 是 C 的两个满足 \(v(xy)=+\infty\) 的元素，则由 \((\mathrm{VL}_{\mathrm{I}})\) 必有 \(v(x)=+\infty\) 或 \(v(y)=+\infty\)；换言之，商环 C/p 除 0 外没有零因子；立即可验证，由 v 通过取商导出的映射 \(\overline{v}:C/p\to\Gamma_{\infty}\) 是 C/p 上的赋值，且这个赋值下 \(+\infty\) 的原像约化为 0。
\end{enumerate}

\subsection{域上的赋值}

命题 2. 设 K 是一个（不必交换的）域，v 是 K 的以 Γ 为值群的赋值。则：

\begin{enumerate}
\def\labelenumi{(\roman{enumi})}
\item
  对 \(x \neq 0\)，\(v(x) \neq +\infty\)。
\item
  集 \(\mathbf{A}\) 由 \(x\in \mathbf{K}\) 中使 \(v(x)\geqslant 0\) 的元素组成，它是 \(\mathbf{K}\) 的子环。
\item
  对所有 \(\alpha \geqslant 0\)（在 \(\Gamma\) 中），集 \(\mathrm{V}_{\alpha}\)（相应地，\(\mathrm{V}_{\alpha}^{\prime}\)）由 \(x \in \mathbf{A}\) 中使 \(v(x) > \alpha\)（相应地，\(v(x) \geqslant \alpha\)）的元素组成，它是 \(\mathbf{A}\) 的双边理想，且每个（左或右）理想 \(\neq (0)\)（\(\mathbf{A}\) 的）都包含某个 \(\mathrm{V}_{\alpha}^{\prime}\)。
\item
  集 \(\mathfrak{m}(\mathbf{A})\) 由 \(x \in \mathbf{A}\) 中使 \(v(x) > 0\) 的元素组成，它是 \(\neq \mathbf{A}\) 的最大理想（在 \(\mathbf{A}\) 中）；\(\mathbf{U}(\mathbf{A}) = \mathbf{A} - \mathfrak{m}(\mathbf{A})\) 是 \(\mathbf{A}\) 的可逆元素所成之集，而 \(\kappa(\mathbf{A}) = \mathbf{A} / \mathfrak{m}(\mathbf{A})\) 是一个（不必交换的）域。
\item
  对所有 \(x \in \mathbf{K} - \mathbf{A}\)，\(x^{-1} \in \mathfrak{m}(\mathbf{A})\)。
\end{enumerate}

断言 (i) 由 \(\overline{v}^{-1}(+\infty)\) 是 K 的一个不等于 K 的理想这一事实得出。A 是环以及诸 \(V_{\alpha}\)、\(V_{\alpha}^{\prime}\) 是双边理想这一事实的验证，凭公理（\(VL_{I}\)）、（\(VL_{II}\)）与（\(VL_{III}\)）是平凡的。若 a 是 A 的一个（例如左）理想，且 \(x \neq 0\) 属于 A，则每个 \(y \in A\)，只要 \(v(y) \geqslant v(x)\)，都可写成 y = zx，其中 \(z = yx^{-1}\)，于是 \(v(z) = v(y) - v(x) \geqslant 0\)，从而 \(z \in A\)；换言之，左理想 Ax 包含诸 \(V_{\alpha}^{\prime}\)（对 \(\alpha \geqslant v(x)\)）。集 \(U(A) = A - m(A)\) 是由 \(x \in K\) 中使 \(v(x) = 0\) 的元素组成之集；若 \(x \in U(A)\)，则

\[
v (x ^ {- 1}) = - v (x) = 0,
\]

从而 \(x^{-1} \in \mathrm{U}(\mathrm{A})\)；反之，若 \(y \in \mathrm{A}\) 在 \(\mathrm{A}\) 中可逆，则 \(v(y) \geqslant 0\)，\(v(y^{-1}) \geqslant 0\)，且 \(v(y) + v(y^{-1}) = 0\)，于是 \(v(y) = 0\)，即 \(y \in \mathrm{U}(\mathrm{A})\)；这证明了 (iv)，而 (v) 由定义立即得出。

A（相应地，m(A)、κ(A)）称为 K 上赋值 v 的环（相应地，理想、剩余域）。

显然 U(A) 是同态 \(v: K^{*} \to \Gamma\) 的核，且像 \(v(\mathbf{K}^{*})\)（乘法群 \(K^{*}\) 在 v 之下的）是加法群 \(\Gamma\) 的子群，称为 v 的阶群或值群，它因此同构于 \(K^{*}/U(A)\)；对 \(x \in K\)，元素 \(v(x)\)（\(\Gamma_{\infty}\) 中的）有时称为 x 关于 v 的赋值或阶。K 上两个赋值 \(v, v'\) 若有相同的环，则称为等价的。

命题 3. 为使两个赋值 \(v, v'\) 在 \(a\)（不必交换的）域 \(K\) 上等价，必须且只需存在同构 \(\lambda\)，它把序群 \(v(K^*)\) 映到序群 \(v'(K^*)\) 上，使 \(v' = \lambda \circ v\)。

设 \(v\) 与 \(v'\) 等价；按假设，赋值 \(v\) 的环 A 与 \(v'\) 的环相同，\(v\) 与 \(v'\)（限制在 K* 上）分解为同态 \(K^* \to K^*/U(A) \xrightarrow{\mu} v(K^*)\) 与 \(K^* \to K^*/U(A) \xrightarrow{\nu} v'(K^*)\)，其中 \(\mu\) 与 \(\nu\) 是同构；此外，\(v(K^*)\)（相应地，\(v'(K^*)\)）的正元素所成之集，是在 \(\mu\)（相应地，\(\nu\)）之下 m(A) 的 \(\neq 0\) 的元素模 U(A) 的类所成之集的像；我们断定 \(\lambda = \nu \circ \overline{\mu}^{-1}\) 解决问题，而逆命题是明显的。

现在设 K 是交换域；这时，对 K 上的每个赋值 v，赋值 v 的环 A 都是 § 1 第 2 节定义 2 意义下 K 的赋值环（这证明了所用术语的合理性）；这由命题 2 (c) 与 § 1 第 2 节定理 1 (c) 立即得出。反之，回忆对每个以 K 为分式域的整环 B，整除关系 x\textbar y（等价于 \(y \in Bx\)）使 \(K^{*}\) 成为一个预序群，其相伴序群 \(\Gamma_{B}\) 是商群 \(K^{*}/U(B)\)，即 \(K^{*}\) 对 B 的可逆元素所成的群 \(U(B)\) 的商，这个群的正元素是 \(B^{*}/U(B)\) 的元素（其中 \(B^{*} = B - \{0\}\)）；映射 \(x \mapsto Bx\) 通过取商定义了从序群 \(K^{*}/U(B)\) 到 K 的非零主分式理想所成的群（以关系 \(\supset\) 排序）上的同构（《代数》第六章 § 1 第 5 节）。以 K 为分式域且使群 \(\Gamma_{A} = K^{*}/U(A)\) 全序的那些环 A 恰是 K 的赋值环（§ 1 第 2 节定理 1 (d)）。若 \(v_{A}\) 表示从 \(K^{*}\) 到 \(\Gamma_{A}\) 上的典范同态，则立即可见 \(v_{A}\)（用 \(v_{A}(0) = +\infty\) 扩张后）是 K 上的一个（称为典范的）赋值，其环为 A；每个等价于 \(v_{A}\) 的赋值都可写成 \(v = \sigma \circ v_{A}\)，其中 \(\sigma\) 是从 \(\Gamma_{A}\) 到 v 取值的群的一个子群上的同构（命题 3）；\(\sigma \circ v_{A}\) 称为 v 的典范分解。

命题 4. 设 C 是整环，K 是它的分式域，\(C^{*} = C - \{0\}\)，\(v: C \to \Gamma_{\infty}\) 是 C 上的赋值。则存在唯一的延拓 v 的 K 上的赋值 w，且 \(w(K^{*})\) 是 \(\Gamma\) 的子群，由 \(v(C^{*})\) 所生成。

由《代数》第一章 § 2 第 7 节定理 2，存在唯一的同态 \(w\)（从 \(\mathbf{K}^*\) 到 \(\Gamma\)），它延拓 \(v \mid \mathbf{C}^*\)，且 \(w(\mathbf{K}^*)\) 由 \(v(\mathbf{C}^*)\) 生成。剩下要证明 \(w\) 满足公理 \((\mathrm{VL}_{\mathrm{II}})\)。于是设 \(x' \in \mathbf{K}^*\)、\(y \in \mathbf{K}^*\) 满足 \(x + y \in \mathbf{K}^*\)；存在 \(a \in \mathbf{C}^*\) 使 \(ax \in \mathbf{C}^*\) 且 \(ay \in \mathbf{C}^*\)，从而 \(a(x + y) \in \mathbf{C}^*\)。由于 \(w\) 在 \(\mathbf{C}^*\) 上的限制满足 \((\mathrm{VL}_{\mathrm{II}})\)，

\[
w (a (x + y)) \geqslant \inf (w (a x), w (a y)).
\]

从两端消去 \(w(a)\)，我们得到

\[
w (x + y) \geqslant \inf (w (x), w (y)).
\]

\subsection{互译}

设 K 是（交换）域，f 是 K 的位，v 是 K 上的赋值，A 是 K 的赋值环。若 A 既是 f 的环又是 v 的环，则称 A、f 与 v 相伴。由第 1 节与 § 2 第 3 节，三个对象 A、f 与 v 中的每一个都确定另外两个（对位与赋值而言，确定到相差一个等价）。特别地，有下列等价关系：

\[
\begin{array}{l l l l} x \in \mathrm{A} & \Leftrightarrow f (x) \neq \infty & \Leftrightarrow v (x) \geqslant 0 \\ x \in \mathfrak {m} (\mathrm{A}) & \Leftrightarrow f (x) = 0 & \Leftrightarrow v (x) > 0 \\ x \in \mathrm{A} - \mathfrak {m} (\mathrm{A}) = \mathrm{U} (\mathrm{A}) \Leftrightarrow f (x) \neq 0 \text {and} f (x) \neq \infty \Leftrightarrow v (x) = 0 \\ x \in \mathrm{K} - \mathrm{A} & \Leftrightarrow f (x) = \infty & \Leftrightarrow v (x) <   0 \end{array}
\]

关于赋值环、位或赋值的每个结果都可以翻译成关于另外两个概念的结果。于是 § 2 第 4 节命题 4 给出：

命题 5. 设 K 是域，v 是 K 上的赋值，K' 是 K 的扩张。存在 K' 上的赋值 v'，它在 K 上的限制等价于 v。

设 \(\Gamma_v\) 与 \(\Gamma_{v'}\) 是 \(v\) 与 \(v'\) 的阶群。由于 \(v'\) 在 K 上的限制等价于 \(v\)，存在同构 \(\lambda\)，它把 \(\Gamma_v\) 映到 \(\Gamma_{v'}\) 的一个子群上，使在 K 上 \(v' = \lambda \circ v\)。若把 \(\Gamma_v\) 等同于 \(\lambda(\Gamma_v)\)（借助 \(\lambda\)），则可见 \(v'\) 延拓 \(v\)。

注意 \(\Gamma_{v'}\) 一般异于 \(\lambda(\Gamma_v)\)，且 \(v'\) 的等价类不必唯一。我们将在 § 8 中回到这一点。

翻译 § 1 第 3 节定理 3（或 § 2 第 5 节命题 6），我们得到：

命题 6. 设 K 是域，A 是 K 的子环，x 是 K 的元素。为使 x 在 A 上整，必须且只需 K 的每个在 A 上为正的赋值都在 x 处为正。

从现在起，我们将把进行与上述类似的翻译的工作一般留给读者。

\subsection{赋值的例子}

§ 1 第 4 节给出的赋值环的例子为我们提供下面例 1 至例 4：

例 (1) 有限域 F 上的每个赋值都是非真的，因为 \(F^{*}\) 的每个元素都是单位根。

例 (2) 若 K 是域 \(K'\) 的子域，则 \(K'\) 上的赋值在 K 上的限制是 K 上的赋值。

例 (3) 设 \(k\) 是域，\(K = k((T))\)。把每个非零形式幂级数映为它的阶（《代数》第四章 § 5 第 7 节）的映射 \(v\) 是 \(K\) 上的赋值，其阶群为 \(Z\)，环为 \(k[[T]]\)。相伴的位是典范同态 \(f: k[[T]] \to k\) 向 \(k((T))\) 的扩张，其中令 \(f(u) = \infty\)（当 \(u \notin k[[T]]\) 时）。

例 (4) 设 A 是主理想整环，K 是它的分式域，p 是 A 的一个极端元。对 \(x \in K^{*}\)，用 \(v_{p}(x)\) 表示 p 在 x 的分解为极端元素的分解式中的指数（《代数》第七章 §1 第 3 节定理 2）；立即可见 \(v_{p}\) 是一个赋值，其阶群为 Z，环为 \(A_{Ap}\)。由 §1 第 4 节命题 3，我们如此得到（在相差一个等价的意义下）K 上的所有非真的且在 A 上为正的赋值。取 A = Z，我们重新得到 Q 上的 p-进赋值（《一般拓扑学》第九章 §3 第 2 节）；这些赋值是（在相差一个等价的意义下）Q 上仅有的非真赋值（§1 第 4 节命题 3 的推论 1）。取 A = k{[}X{]}（其中 k 是域），k(X) 上的限制在 k 上为非真的那些非真赋值是（在相差一个等价的意义下）：一方面是诸赋值 \(v_{P}\)，其中 P 取遍 k{[}X{]} 的首一不可约多项式的集合；另一方面是由

\[
v (\mathrm{P} / \mathrm{Q}) = \deg (\mathrm{Q}) - \deg (\mathrm{P})
\]

（其中 \(P \in k[X]\)、\(Q \in k[X]\)）所定义的赋值 v（§ 1 第 4 节命题 3 的推论 2）；所有这些赋值显然都以 Z 为阶群，其剩余域是 k 的单生成代数扩张（《代数》第五章 § 3 第 1 节）。

例 (5) 映射 \(P(X, Y) \mapsto P(T, e^{\mathrm{T}})\)（从 \(\mathbf{C}[X, Y]\) 到 \(\mathbf{C}((T))\)）是单射（《实变函数》第四章 §2 命题 9），因此可以扩张为从 \(\mathbf{C}(X, Y)\) 到 \(\mathbf{C}((T))\) 的一个子域上的同构。例 3 中定义的 \(\mathbf{C}((T))\) 上的赋值在这个子域上的限制定义了 \(\mathbf{C}(X, Y)\) 上的一个赋值，它在 C 上非真，其阶群为 Z，剩余域为 C。

第 2 节命题 4 使我们能够构造阶群与剩余域都给定的赋值：

例 (6) 设 \(\Gamma\) 是全序群，\(k\) 是域。设 \(\Gamma_{+}\) 是 \(\Gamma\) 的正元素所成的幺半群，C 是 \(\Gamma_{+}\) 在 \(k\) 上的代数。按定义，C 有一组基 \((x_{\alpha})_{\alpha \in \Gamma_{+}}\)（在 \(k\) 上），其乘法表为 \(x_{\alpha}x_{\beta} = x_{\alpha +\beta}\)。若 \(x = \sum_{\alpha}a_{\alpha}x_{\alpha}\) 是 C 的非零元素，记 \(v(x) = \inf_{a_{\alpha}\neq 0}(\alpha)\)，并记 \(v(0) = +\infty\)；立即可验证映射 \(v\)（从 C 到 \(\Gamma_{\infty}\)）满足第 1 节的条件 \((\mathrm{VL}_{\mathrm{I}})\) 与 \((\mathrm{VL}_{\mathrm{II}})\)，且 C 是整环。设 K 是 C 的分式域，\(w\) 是 K 上延拓 \(v\) 的赋值（第 2 节命题 4）。由于 \(\Gamma\) 的每个元素都是两个正元素之差，\(w\) 以 \(\Gamma\) 为阶群。设 A 是 \(w\) 的环，m 是它的极大理想；我们将证明 A 是 m 与 k（等同于 k.1）的直和，这将证明 \(w\) 的剩余域同构于 k。显然 \(m \cap k = (0)\)。另一方面，用 p 表示 C 的由诸 \(x_{\alpha}\)（其中 \(\alpha > 0\)）生成的理想，K 中每个赋值为 0 的元素 x 都可写成 \((a + y)/(b + z)\) 的形式，其中 \(a \in k^{*}\)、\(b \in k^{*}\)、\(y \in p\)、\(z \in p\)；则

\[
x = a b ^ {- 1} + (b y - a z) b ^ {- 1} (b + z) ^ {- 1}
\]

从而 \(w(x - ab^{-1}) > 0\)，且 \(x \equiv ab^{-1} \pmod{m}\)；这就表明了我们的断言。

若 \(\Gamma = \mathbf{Z} \times \mathbf{Z}\)，则 \(\mathbf{K} = k(\mathbf{X}, \mathbf{Y})\)，于是上述构造给出 \(k(\mathbf{X}, \mathbf{Y})\) 上的一批在 \(k\) 上非真的赋值，其阶群为 \(\mathbf{Z} \times \mathbf{Z}\)，剩余域为 \(k\)。这些赋值依赖于 \(\mathbf{Z} \times \mathbf{Z}\) 上所选的序结构。例如，可以给 \(\mathbf{Z} \times \mathbf{Z}\) 赋以字典序。或者，对无理数 \(\alpha\)，\(\mathbf{Z} \times \mathbf{Z}\) 可等同于 \(\mathbf{R}\) 的一个子群（借助同态 \((m, n) \mapsto m + n\alpha\)，由于 \(\alpha\) 是无理数，这个同态是单射），并赋以由 \(\mathbf{R}\) 上的序诱导的序。

利用第 2 节命题 4 的赋值的其他构造将在 § 10 中描述。

\subsection{赋值环的理想}

定义 2. 设 G 是有序集。若 G 的子集 M 满足：关系 \(x \in \mathbf{M}\) 与 \(y \geqslant x\) 蕴涵 \(y \in \mathbf{M}\)，则称 M 为优势集。

设 K 是域，v 是 K 上的赋值，A 是 v 的环，G 是 v 的阶群。对每个优势子集 \(M \subset G\)，用 \(\mathfrak{a}(M)\) 表示由 \(x \in K\) 中使 \(v(x) \in M \cup \{+\infty\}\) 的元素组成之集。显然 \(\mathfrak{a}(M)\) 是 K 的子 A-模。

命题 7. 映射 \(\mathbf{M} \mapsto \mathfrak{a}(\mathbf{M})\) 是从 \(\mathbf{G}\) 的优势子集所成之集到 \(\mathbf{K}\) 的子 A-模所成之集上的递增双射。

设 \(\mathfrak{b}\) 是 \(\mathbf{K}\) 的子 A-模。由诸 \(v(x)\)（\(x \in \mathfrak{b} - (0)\)）组成的集是优势子集 \(\mathbf{M}(\mathfrak{b})\)（\(\mathbf{G}\) 的）。若能证明下列等式，命题 7 即得证：

\begin{enumerate}
\def\labelenumi{(\arabic{enumi})}
\setcounter{enumi}{1}
\item
  \(\mathbf{M}(\mathfrak{a}(\mathbf{N})) = \mathbf{N}\)（对每个优势子集 \(\mathbf{N}\)，\(\mathbf{G}\) 中的）；
\item
  \(\mathfrak{a}(\mathbf{M}(\mathfrak{b})) = \mathfrak{b}\)（对每个子 A-模 \(\mathfrak{b}\)，\(\mathbf{K}\) 的）。
\end{enumerate}

公式 (2) 是容易的，因为对所有 \(m \in \mathbf{N}\)，存在 \(x \in \mathbf{K}\) 使 \(v(x) = m\)。于是显然 \(\mathfrak{b} \subset \mathfrak{a}(\mathbf{M}(\mathfrak{b}))\)；反之，设 \(x \in \mathfrak{a}(\mathbf{M}(\mathfrak{b}))\)，并设 \(x \neq 0\)；

则 \(v(x) \in \mathbf{M}(\mathfrak{b})\)，从而存在 \(y \in \mathfrak{b}\) 使 \(v(x) = v(y)\)；于是 \(x = uy\)，其中 \(v(u) = 0\)，这证明 \(x \in \mathrm{Ay} \subset \mathfrak{b}\)，证明完成。

推论. 设 \(G_{+}\) 是 G 中正元素所成之集。映射 \(\mathbf{M} \mapsto \mathfrak{a}(\mathbf{M})\) 是从 \(G_{+}\) 的优势子集所成之集到 A 的理想所成之集上的双射。

由于 \(\mathbf{A} = \mathfrak{a}(\mathbf{G}_{+}),\mathfrak{a}(\mathbf{M})\subset \mathbf{A}\) 等价于 \(\mathbf{M}\subset \mathbf{G}_{+}\)。

例如，极大理想 \(m(A)\) 等于 \(a(S)\)，其中 S 表示 G 的严格正元素所成之集。
